% ---------------------------------------------
% Copyright Ben Paul Wise. All rights reserved.
% ---------------------------------------------
% process with TeXworks as pdfLaTeX
%
% ACV --- Derivation of Record (P1, round 3, 2026-07-25).
% Reconstructed 2026-10-09 from ACV-derivation-of-record.pdf, whose LaTeX
% source was lost.  The text, the equation identifiers (E-n) and the tables
% follow that PDF; the page format comes from the shared prolog.

\input{../../CPVI/shared-tex/prolog-article.tex}

\usepackage{amsthm}
\usepackage{booktabs}
\usepackage{array}
\usepackage{longtable}

\def\myversion{P1 r3}
\setlength{\headheight}{15pt}

\lfoot{\color{myDarkRed}\bf{DRAFT \myversion}}
\cfoot{{\thepage} of {\pageref*{LastPage}}} % asterix suppressed hyperlink
\rfoot{\copyright \  Ben P. Wise}

% Propositions, corollaries, definitions and remarks share one counter,
% numbered within the section; the two assumptions are numbered on their own.
\theoremstyle{plain}
\newtheorem{proposition}{Proposition}[section]
\newtheorem{corollary}[proposition]{Corollary}
\theoremstyle{definition}
\newtheorem{definition}[proposition]{Definition}
\newtheorem{assumption}{Assumption}
\theoremstyle{remark}
\newtheorem{remark}[proposition]{Remark}

% Plain ruled box for the statements the document singles out.
\newtcolorbox{recordbox}{enhanced, breakable, sharp corners,
  colback=white, colframe=black, boxrule=0.6pt,
  left=4pt, right=4pt, top=2pt, bottom=2pt,
  before upper={\setlength{\parskip}{0.6em}}}

% Part headings: Part A, Part B, Part C.
\newcommand{\docpart}[1]{\par\vspace{3ex}\noindent
  {\color{myBlue1}\huge\bfseries #1}\par\vspace{1ex}
  \addcontentsline{toc}{part}{#1}}

\newcommand{\Sc}{\mathcal{S}}
\newcommand{\Zc}{\mathcal{Z}}
\newcommand{\bbA}{\mathbb{A}}
\newcommand{\E}{\mathbb{E}}
\newcommand{\ch}{\mathrm{ch}}
\newcommand{\eff}{\mathrm{eff}}

% Equation identifiers.  Single-line equations carry \tag{E-n}.  A stacked
% group is an eqnarray whose rows are E-nA, E-nB, ...; \Elabel{E-1A}{E1_A}
% at the end of a row sets the number printed for that row and labels it.
\makeatletter
\newcommand{\Etagname}{}
\renewcommand{\theequation}{\Etagname}
\newcommand{\Elabel}[2]{\gdef\Etagname{#1}\phantomsection
  \protected@edef\@currentlabel{#1}\label{#2}}
\makeatother
% \tag'd equations do not advance the equation counter, so hyperref's default
% anchor names repeat and links land on the wrong equation.  Unique names:
\hypersetup{hypertexnames=false}

\begin{document}

\title{ACV --- Derivation of Record\\[1ex]
  {\large Inferring which entries of a transition matrix are exactly zero,
   without assuming anything about where they are}\\[1ex]
  {\normalsize v1: $K=0$; one entrywise sparsity model (D-30)}}
\author{Agent 1 (Derivation Author)}
\date{P1, round 3 --- July 25, 2026}

\maketitle


\section*{What this is}

You have $N$ time series measured together --- prices, flows, temperatures, counts
of something --- sampled at regular intervals. You believe each series is driven
partly by its own past and partly by the past of the others, and you want to know
\emph{which} influences are real. Write the belief down as a linear transition rule:
today's vector of $N$ numbers is a matrix times yesterday's vector, plus noise. That
matrix has $N^2$ entries; with $p$ lags of history it has $N^2p$. At $N=40$ and $p=2$
that is 3200 numbers, fitted from perhaps two hundred observations. Ordinary least
squares will return 3200 non-zero numbers, most of them noise, and no indication of
which ones matter.

\textbf{What this method does.} You hand it the series. It returns a transition
matrix in which most entries are \emph{exactly zero} --- not small, zero ---
together with a posterior probability for each entry being non-zero, and a posterior
over how many non-zero entries there are. \textbf{It decides how many by weighing
evidence, not by anyone's judgement about which variables ought to influence which.}
There is no significance test, no cross-validation, no threshold on coefficient size,
and no information criterion anywhere in it.

\subsection*{The core formula}

Two factors, multiplied. Let $M$ be the number of entries of the transition matrix
that are non-zero, let $x_{ti}$ be the observed value of coordinate $i$ at time $t$,
and let $u_{ti}$ be the value the matrix multiplication predicts for it from the
previous observations. Then the object being maximised is, in its recognisable form,
\begin{equation}
\underbrace{\;\;p^{M}\;\;}_{\text{prior mass on the pattern}}
\;\times\;
\underbrace{\prod_{t}\prod_{i}\exp\Bigl\{-\tfrac12\Bigl(\frac{x_{ti}-u_{ti}}{\sigma_i}\Bigr)^{2}\Bigr\}}_{\text{Gaussian on the residual}} .
\tag{$\star$}\label{star}
\end{equation}
The left factor is the prior: each entry is switched on independently with
probability $p$, so a pattern with $M$ of them switched on carries prior mass
proportional to $p^M$. The right factor is the likelihood: the residual between what
was observed and what the matrix predicts, measured in units of that coordinate's
noise scale $\sigma_i$, and taken as Gaussian. Nothing more than this is needed to
state what the method is.

\subsection*{What balancing the two factors means}

Switching on more entries makes $u_{ti}$ track $x_{ti}$ more closely, so the right
factor rises. It also makes $M$ larger, so the left factor falls, since $p<1$.
Switching off entries does the opposite: the prior mass rises and the fit degrades.
\textbf{The posterior picks the trade.} There is a rate of exchange between fit and
prior mass, and where it turns is where the answer is. That balance is what is meant
by an Occam factor in the rest of this document, and \textbf{it is the entire
complexity-control mechanism of the project} --- everything else is either an exact
evaluation of one of the two factors, a statement of what the evaluation assumes, or
a case where it degenerates.

\begin{recordbox}
\textbf{The essential property, and it is the point of the project.}

\textbf{The method assumes nothing whatsoever about \emph{where} the non-zero entries
lie.} There is no notion of a block, a group, a partition, a neighbourhood, a
distance, an ordering, or a position anywhere in the model. Every one of the $N^2p$
entries is exchangeable with every other, a priori. Look at \eqref{star}: the prior
factor is $p^M$, and $M$ is a \emph{count}. There is nothing in that expression that
can see which entries are on --- only how many.

It follows that if a fitted result comes out looking block-structured, or banded, or
clustered by country, or clustered by anything at all, \textbf{that is a finding
about the data and never an input to the model.} The machinery has no vocabulary in
which to have preferred it.

This is why the method exists. GVAR --- the standard tool for the motivating problem
--- has its designers write the link structure down in advance from economic
intuition about which countries trade with which. Not doing that is the whole
contribution. It is emphatically \textbf{not} a goal to reproduce GVAR's particular
structure; a result that disagreed with GVAR would be just as much a success, and a
method tuned until it agreed would be a failure that had learned to hide.
\end{recordbox}

\subsection*{Lineage, in one paragraph}

AUTOCLASS is a Bayesian classification engine from the late 1980s that decides how
many classes a dataset contains by marginal likelihood alone --- no elbow plots, no
stopping parameter. It ran unmodified on stellar spectra and on satellite imagery.
This project transplants that discipline from \emph{classes over cases} to
\emph{zeros over matrix entries}: same refusal to introduce a criterion outside the
marginal likelihood, same insistence that the complexity penalty fall out of the
prior rather than be imposed. That is the extent of the analogy and it is offered as
motivation, not as structure. Where AUTOCLASS is cited below it is cited from
\texttt{workflow/autoclass-reference.md}, which quotes the primary sources; the
project has twice mis-stated AUTOCLASS from memory and will not do so again.

\subsection*{How to read this, and where the developed formula is}

\textbf{Forward pointer.} \eqref{star} is the shape of the thing, not the thing that
is computed. The fully developed posterior --- the object the implementation
evaluates and the search maximises --- is \eqref{E126} \textbf{in
\S\ref{sec:search}}, assembled from the closed-form per-equation marginal
\eqref{E101}, the structure prior \eqref{E135}, and the one-dimensional integral
\eqref{E103}. It differs from \eqref{star} in exactly three respects --- Gaussian
normalising constants, a prior over the \emph{values} of the non-zero entries, and
the marginalisation of the rate $p$ --- and \textbf{each of the three is introduced
later, at the point where it is needed, with the argument for why it is forced}: the
normalising constants in \S\ref{sec:model}, the prior over the values in
\S\ref{sec:slab}, and the marginalisation of the rate in \S\ref{sec:rate}. Gaussian
errors and independent matrix entries are used exactly as \eqref{star} suggests;
nothing is quietly replaced.

\begin{description}[style=nextline, leftmargin=2em, topsep=0pt, itemsep=0pt]
\item[Part A, \S1--\S4: everything, at one lag.]
  The model, the sufficient statistics, the prior with every constant it contains,
  the closed-form marginal likelihood, and the resulting rule for when an entry gets
  switched on. A reader who stops at the end of Part A has seen every mechanism in
  the project. \emph{Read this part.}
\item[Part B, \S5: the worked example.]
  Benchmark B1 --- fifty synthetic coordinates, 350 true links out of 2500,
  deliberately non-economic --- with the detection threshold and the expected
  false-positive count computed in advance, so that the run can refute the theory
  rather than merely illustrate it. \emph{Read the table.}
\item[Part C, \S6--\S11: the rest.]
  General lag order; the search; the degenerate cases; three reductions; the one
  place the model class is provably inadequate (over-differencing, \S\ref{sec:cannot}); a
  conditioning diagnostic; assumptions and limitations. \emph{Skippable on a first
  pass except \S\ref{sec:red3}, which is the correctness check with teeth.}
\end{description}
Notation is introduced where it is used. Appendix \ref{app:disp} is the
equation-identifier disposition, including what this revision deleted.

\subsection*{What changed in this revision, and why}

\textbf{D-30, ruled 2026-07-25.} The previous draft carried \emph{two} sparsity
models: the entrywise one above, and a second in which the coordinates were
partitioned into groups and inclusion indicators were placed at the block level, with
the partition inferred. The second model, the group count, the block indicators, the
comparison between the two models and every claim about which should win on which
data are \textbf{deleted}. They are preserved in
\texttt{theory/theory-v1-with-blocks.tex} and their identifiers
\textbf{E-114--E-117} are retired permanently.

The reason is the boxed property above. Block structure does not need a block prior
to be recovered: entrywise inference is indifferent to position, so it finds
non-zeros wherever they lie, clustered or scattered. A block prior buys only
strength-borrowing within a block. What it cost while it stood was worse: it could
not express a partly-filled block, so on a benchmark whose truth was 350 entries it
was forced to report 500; and it invited the entirely reasonable suspicion that the
method \emph{assumed} the answer. It never did. But a reader could not tell that from
the derivation, and a derivation from which the central property cannot be read is
defective whatever its algebra says.

Three corrections were also made, and one of them changes a headline number.
\begin{enumerate}[label=(\roman*), topsep=0pt, itemsep=0pt]
\item \textbf{The false-positive prediction was wrong by roughly a factor of
  $(\bar k+1)$} --- it was evaluated as though the retained structure were empty
  when it in fact carries about seven entries per equation. Corrected at
  \eqref{E120}, \eqref{E122}, \eqref{E133} and throughout \S\ref{sec:bench}. The
  corrected result \textbf{carries the mean sparsity $\bar k$ and is therefore not
  parameter-free.} It was described as parameter-free and that was wrong; see \S5.3.
\item \textbf{The claim that no quantity needs a user-supplied value was false.}
  Section \ref{sec:constants} now declares every set constant openly, with reasoning,
  in one table --- including an evidence threshold that the quadrature grid and the
  stopping rule are both chained to and that nobody had named.
\item \textbf{A third reduction with teeth} was added (\S\ref{sec:red3}). The
  existing two are both evaluated at points where every mechanism under test is
  switched off.
\end{enumerate}


\docpart{Part A --- The complete inference at one lag}

\section{The model, and what the data reduces to}\label{sec:model}

\subsection{The process}

Observed $y_t\in\mathbb{R}^N$ for $t=1,\dots,T$; differencing of order $d$ has
already been applied and $y_t$ denotes the differenced series ($d=1$ fixed, D-05 ---
\S\ref{sec:cannot} is what that costs). There is no hidden state (A-001), so the
observed vector is the state:
\begin{eqnarray}
y_t &=& \sum_{l=1}^{p} A_l\, y_{t-l} \;+\; c \;+\; \varepsilon_t,
  \Elabel{E-1A}{E1_A} \\
\varepsilon_t &\overset{\text{iid}}{\sim}& \mathcal{N}(0,\Sigma),
  \Elabel{E-1B}{E1_B} \\
A_l &\in& \mathbb{R}^{N\times N},
  \Elabel{E-1C}{E1_C}
\end{eqnarray}
with $c\in\mathbb{R}^N$ an intercept. The observation equation \textbf{E-2} is the
identity, $y_t=Jx_t$ with $J=I_N$, and carries nothing. Part A takes $p=1$;
\S\ref{sec:lag} carries general $p$, and every statement here is written so that the
generalisation is a substitution.

\paragraph{The intercept is always included and is never a sparsity candidate.} The
previous draft had no intercept, which on differenced data asserts zero drift for
every coordinate. Ben ruled it in. It costs $N$ parameters, changes no derivation,
and is absorbed by writing the regressor vector with a leading constant:
\begin{eqnarray}
z_t &=& \bigl(1,\; y_{t-1}',\; \dots,\; y_{t-p}'\bigr)' \;\in\; \mathbb{R}^{1+Np},
  \Elabel{E-134A}{E134_A} \\
\bbA &=& \bigl[\, c \;\; A_1 \;\cdots\; A_p \,\bigr],
  \Elabel{E-134B}{E134_B} \\
y_t &=& \bbA z_t + \varepsilon_t.
  \Elabel{E-134C}{E134_C}
\end{eqnarray}
Index 0 of $z_t$ is the constant. Every $\gamma_i$ below contains index 0 by
construction; the candidate set is indices $1,\dots,Np$, so the candidate count per
equation is $n_i=Np$ and the global count is $n=N^2p$, exactly as before. Two things
this buys: the $p=0$ reduction now lands on AUTOCLASS's real-valued model as shipped
rather than on that model with its mean pinned at zero (\S\ref{sec:red2}), and a
coordinate with genuine drift no longer forces spurious diagonal links to absorb it.

The innovation covariance is diagonal (D-04),
\begin{eqnarray}
\Sigma &=& \operatorname{diag}(\sigma_1^2,\dots,\sigma_N^2),
  \Elabel{E-4A}{E4_A} \\
\sigma_i^2 &>& 0,
  \Elabel{E-4B}{E4_B}
\end{eqnarray}
and the parameter set and its dimension are
\begin{eqnarray}
\Theta &=& \bigl\{(c, A_1,\dots,A_p,\sigma_1^2,\dots,\sigma_N^2)\bigr\},
  \Elabel{E-5A}{E5_A} \\
D &=& N^2p+2N,
  \Elabel{E-5B}{E5_B}
\end{eqnarray}
before any entry is forced to zero. Under a structure that zeroes some entries of the
$A_l$, $D$ falls by exactly the number excluded. \textbf{That count is a property of
the model and not of a choice of coordinates}, because with no hidden state there is
no basis to change --- which is why ``how many entries are non-zero'' is a well-posed
question here and was not in the $K>0$ draft.

\begin{assumption}[conditioning, and the whole of it]\label{ass:cond}
Every marginal likelihood in this document is conditional on $y_{1-p_{\max}:0}$, with
$p_{\max}$ the largest lag order in the candidate set, \textbf{fixed before any
fitting}. Every candidate therefore explains the identical vector $y_{1:T}$ from the
identical number of usable transitions. $p_{\max}$ is a declared quantity, listed in
\S\ref{sec:constants}; it is \emph{not} the reported ceiling $p^*$ of
\S\ref{sec:ceiling}, and conflating the two was a defect of the previous draft.
\end{assumption}

With that, the likelihood is the complete-data likelihood --- there is nothing latent
to integrate:
\begin{equation}
f\bigl(y_{1:T}\mid y_{1-p_{\max}:0},\theta,\Sc\bigr) \;=\;
\prod_{t=1}^{T}\mathcal{N}\bigl(y_t;\; \bbA z_t,\; \Sigma\bigr),
\tag{E-9}\label{E9}
\end{equation}
which is the complete-data likelihood \textbf{E-10} of the earlier drafts evaluated
at $x=y$ --- with no hidden state the complete-data and observed-data likelihoods are
the same object. Write $\ell(\theta)$ for its logarithm. \textbf{No Kalman recursion,
no EM, no Cheeseman--Stutz substitution, and no Laplace approximation is used to
compute any reported number in v1}; \S\ref{sec:marginal} evaluates the marginal
likelihood exactly up to one one-dimensional integral.

\paragraph{First refinement: the Gaussian normalising constants.} The right factor of
\eqref{star} is written as a bare exponential, which is how one writes it when only
the residuals are of interest. \eqref{E9} is the full density, so it carries a factor
$(2\pi\sigma_i^2)^{-1/2}$ for each observation of each coordinate. The reason those
factors cannot be dropped here is that $\sigma_i$ \textbf{is inferred, not fixed}. If
$\sigma_i$ were known, the constants would be the same for every candidate structure
and would cancel from every comparison, which is why a reader accustomed to
residual-sum-of-squares fitting rarely sees them. Because $\sigma_i$ is a parameter
to be integrated out, the constants are a function of it, and dropping them changes
the shape of the integrand in $\sigma_i^2$: the likelihood would then be monotone
decreasing in $\sigma_i^2$ with no interior maximum, and the inferred dispersion
would be biased towards zero without bound. They cost nothing --- they are a term
$-\frac{T}{2}\log(2\pi\sigma_i^2)$ per equation in \eqref{E28_A} --- and they are what
makes the $\sigma^2$ integral in Proposition \ref{prop:closed} a proper
inverse-Gamma rather than a divergence.

\begin{remark}[the prior is not truncated to stationarity]\label{rem:trunc}
A truncation to $\{\rho(A)<1\}$ has a normalising constant that depends on the
inclusion pattern, and that constant would enter every model comparison --- exactly
the smuggled quantity D-10 forbids. \eqref{E9} is a proper density in $y_{1:T}$ for
every $\theta$, explosive or not, so no truncation is needed for integrability. The
cost is prior mass on explosive processes, recorded in \S\ref{sec:assump}.
\end{remark}

\subsection{What the data reduces to}

\begin{definition}[the statistics]\label{def:stats}
\begin{eqnarray}
S_{xx} &=& \sum_{t=1}^{T} y_t y_t',
  \Elabel{E-26A}{E26_A} \\
S_{xz} &=& \sum_{t=1}^{T} y_t z_t',
  \Elabel{E-26B}{E26_B} \\
S_{zz} &=& \sum_{t=1}^{T} z_t z_t',
  \Elabel{E-26C}{E26_C}
\end{eqnarray}
together with $T$. Their total dimension is $O(N^2p^2)$ and is \textbf{independent
of} $T$.
\end{definition}

\begin{proposition}[sufficiency]\label{prop:suff}
$(T,S_{xx},S_{xz},S_{zz})$ is sufficient for $\theta$ in the model \eqref{E9}.
\end{proposition}

\begin{proof}
Expanding the exponent of \eqref{E9},
\begin{equation}
\ell(\theta) = -\tfrac{TN}{2}\log 2\pi - \tfrac{T}{2}\log|\Sigma|
 - \tfrac12\operatorname{tr}\Bigl\{\Sigma^{-1}\bigl(S_{xx} - \bbA S_{xz}'
 - S_{xz}\bbA' + \bbA S_{zz}\bbA'\bigr)\Bigr\},
\tag{E-27}\label{E27}
\end{equation}
using $u'Bv=\operatorname{tr}(Bvu')$ three times. \eqref{E27} depends on the data
only through $(T,S_{xx},S_{xz},S_{zz})$; by Fisher--Neyman, with the $\theta$-free
factor equal to 1, these are sufficient.
\end{proof}

\begin{remark}[sufficiency does not need a diagonal $\Sigma$, and this is not a
technicality]\label{rem:diag}
\eqref{E27} is derived with $\Sigma$ an \textbf{arbitrary symmetric positive definite
matrix}. Diagonality is used nowhere in Proposition \ref{prop:suff}. It is used for
the first time in \eqref{E28_A} below, to \emph{separate} an object that is already in
hand.

The objection this forestalls is the natural one: \emph{``the sufficient statistics
were derived under a diagonal $\Sigma$, and the per-equation decoupling is then read
off them, so the decoupling argument is circular.''} It is not. Sufficiency is
established for the general-$\Sigma$ model; diagonality is an additional hypothesis
applied afterwards, and the order matters. Agent 3 verified \eqref{E27} symbolically
at $N=2$, $p=1$, $T=2$ with a full symmetric $\Sigma$ precisely so that this rests on
a check rather than on an assertion, and the check must be re-run at $p=2$ in the
same form.
\end{remark}

\begin{corollary}[prior-agnostic consequence]\label{cor:agn}
For any \emph{proper prior}, the marginal likelihood is a function of
$(T,S_{xx},S_{xz},S_{zz})$ and of the structure alone. No candidate structure
requires a pass over the data: the one-time accumulation of \eqref{E26_A}--\eqref{E26_C} costs
$O(TN^2p^2)$ and every subsequent score is $O(k^3)$ or less.
\end{corollary}

\paragraph{Decoupling.} With $\Sigma$ diagonal and $\alpha_i$ the $i$-th row of $\bbA$
written as a column, \eqref{E27} separates:
\begin{eqnarray}
\ell(\theta) &=& \sum_{i=1}^{N}\Bigl[-\tfrac{T}{2}\log(2\pi\sigma_i^2)
 - \tfrac{1}{2\sigma_i^2}\bigl(c_i - 2\,\alpha_i' s_i + \alpha_i' S_{zz}\alpha_i\bigr)\Bigr],
  \Elabel{E-28A}{E28_A} \\
c_i &:=& [S_{xx}]_{ii},
  \Elabel{E-28B}{E28_B} \\
s_i &:=& ([S_{xz}]_{i\cdot})'.
  \Elabel{E-28C}{E28_C}
\end{eqnarray}
Restricting to a structure $\gamma_i\subseteq\{0,1,\dots,Np\}$ (always containing 0)
deletes the corresponding rows and columns:
\begin{eqnarray}
\ell_i\bigl(\alpha_i[\gamma_i],\sigma_i^2\bigr) &=& -\tfrac{T}{2}\log(2\pi\sigma_i^2)
 - \tfrac{1}{2\sigma_i^2}\Bigl(c_i - 2\,\alpha_i[\gamma_i]' s_i[\gamma_i]
 + \alpha_i[\gamma_i]' S_\gamma\, \alpha_i[\gamma_i]\Bigr),
  \Elabel{E-29A}{E29_A} \\
S_\gamma &:=& S_{zz}[\gamma_i,\gamma_i].
  \Elabel{E-29B}{E29_B}
\end{eqnarray}
Equation $i$ needs only $(T,c_i,s_i,S_{zz})$, and $S_{zz}$ is \emph{shared} by all
$N$ equations. At $N=50$, $p=2$ the whole problem is about 79 KiB of statistics at
any $T$.

\begin{remark}[what decouples and what does not]\label{rem:decouple}
\eqref{E28_A} separates the log-\emph{likelihood} across equations. The \emph{score}
separates only conditional on the shared hyperparameters. Under the priors of
\S\ref{sec:prior} there is exactly one shared hyperparameter left --- the slab scale
$g$, which D-09/D-24 require be integrated rather than fixed --- so the correct
statement is: \textbf{$N$ independent per-equation searches, inside one shared
one-dimensional quadrature.} Nothing else couples them. The previous draft also
coupled them through a single global inclusion rate; \S\ref{sec:rate} carries the
per-equation rate instead, which removes that coupling and has a consequence for the
false-positive count that \S\ref{sec:predict} states plainly.
\end{remark}

\subsection{The cost of moving one entry}

This is the result the entire search cost rests on, so it is proved rather than
quoted.

\begin{proposition}[cost of one indicator flip]\label{prop:flip}
\textbf{[Requires: a fixed value of $g$; unconditionally, multiply by the node count
$n_Q$]} Let equation $i$ have structure $\gamma_i$, $k=|\gamma_i|$, with the Cholesky
factor of $\Lambda_\gamma := S_\gamma + g^{-1}I_k$ in hand together with
$\log\det\Lambda_\gamma$ and the scalar
$s_i[\gamma_i]'\Lambda_\gamma^{-1}s_i[\gamma_i]$. Then the score of a neighbour
$\gamma_i\cup\{j\}$ or $\gamma_i\setminus\{j\}$ costs
\begin{equation}
O(k^2)\text{ floating-point operations, independent of $T$ and of $N$,}
\tag{E-85}\label{E85}
\end{equation}
and the score of the whole structure follows in a further $O(1)$.
\end{proposition}

\begin{proof}
Adding regressor $j$ borders the matrix. With $R^{\ch}_\gamma$ the upper Cholesky
factor of $\Lambda_\gamma$,
\begin{eqnarray}
R^{\ch}_{\gamma\cup j} &=&
  \begin{bmatrix} R^{\ch}_\gamma & r \\ 0 & \varrho \end{bmatrix},
  \Elabel{E-30A}{E30_A} \\
(R^{\ch}_\gamma)' r &=& \Lambda[\gamma,j],
  \Elabel{E-30B}{E30_B} \\
\varrho &=& \bigl(\Lambda[j,j] - r'r\bigr)^{1/2},
  \Elabel{E-30C}{E30_C}
\end{eqnarray}
one triangular solve of size $k$, i.e.\ $k^2/2+O(k)$ operations, giving the new
log-determinant as $2(\log\det R^{\ch}_\gamma + \log\varrho)$. The new quadratic form
follows from the rank-one bordering identity
\begin{equation}
s_{\cup j}'\Lambda_{\cup j}^{-1}s_{\cup j} \;=\; s_\gamma'\Lambda_\gamma^{-1}s_\gamma
\;+\; \frac{\bigl(s_{i,j} - r'\,(R^{\ch}_\gamma)^{-\prime} s_\gamma\bigr)^2}{\varrho^2},
\tag{E-86}\label{E86}
\end{equation}
at $O(k)$ cost from quantities already computed in \eqref{E30_B}--\eqref{E30_C}. Deletion is the same
identity in reverse after a Givens downdate, again $O(k^2)$. $T$ appears nowhere: it
entered once, in the accumulation of \eqref{E26_A}--\eqref{E26_C}. $N$ appears nowhere: the other
equations' factors are unchanged and held in log space. The prior term is a ratio of
Beta functions in $k$ alone and is $O(1)$ by \eqref{E113}.
\end{proof}

\begin{corollary}[what this buys]\label{cor:buys}
A sweep over all $n_i=Np$ single-flip neighbours of one equation costs $O(n_ik^2)$
per quadrature node. At $N=50$, $p=2$, $k\approx 7$, $n_Q=32$: about
$1.6\times10^5$ operations per equation-sweep and $8\times10^6$ for all fifty ---
\textbf{independent of} $T$, so the $T=5000$ end of the benchmark sweep costs the
same per move as the $T=20$ end.
\end{corollary}


\section{The prior, and every constant in it}\label{sec:prior}

This is where the audit for smuggled intuition happens. Every constant that enters
the model enters in this section, in view, with a reason --- or it does not enter at
all. \S\ref{sec:constants} is the complete list, including the ones that are
\emph{not} in the model but nonetheless get set.

\subsection{The innovation variances}

Conjugate to \eqref{E29_A}, independently across $i$:
\begin{eqnarray}
\sigma_i^2 &\sim& \operatorname{InvGamma}(a_\sigma,b_\sigma),
  \Elabel{E-92A}{E92_A} \\
a_\sigma = b_\sigma &=& \tfrac12.
  \Elabel{E-92B}{E92_B}
\end{eqnarray}
Read as pseudo-counts, $a_\sigma$ is half the number of prior observations and
$b_\sigma$ half their sum of squares, so \eqref{E92_B} is worth one observation of unit
variance. Coordinates are standardized before fitting (charter \S3), which is what
makes ``unit'' mean something; $b_\sigma/a_\sigma=1$ then follows, and properness
forbids $a_\sigma\to0$, which is Jeffreys. Two constraints, one free scalar left, and
it is set to the value that makes the prior worth exactly one observation.
\textbf{This is an argument for the ratio and for properness; it does not by itself
select $\frac12$ over $\frac14$}, and \S\ref{sec:constants} records it as set rather
than derived, with its magnitude.

\begin{assumption}[what a prior must be]\label{ass:proper}
Every $\pi(\theta\mid\Sc)$ is a \emph{proper} density and $\pi(\Sc)$ a proper mass
function. This is a correctness requirement, not hygiene: an improper prior with an
arbitrary scale $c$ on each of $D_\Sc$ coordinates multiplies the marginal likelihood
by $c^{D_\Sc}$, and since $D_\Sc$ differs across structures by one per entry, every
Bayes factor is then wrong by $(D_{\Sc_1}-D_{\Sc_2})\log c$ --- in the direction that
empties the model as $c\to\infty$. The reflexes a reader brings from the BVAR
literature --- Jeffreys' $1/\sigma^2$, a flat prior on $A$ --- are both excluded.
\end{assumption}

\subsection[Second refinement: the slab]{Second refinement: a prior over the values of the non-zero
entries}\label{sec:slab}

\eqref{star} puts a prior on \emph{which} entries are non-zero and says nothing about
\emph{what they are}. That gap has to be closed before anything can be computed,
because the quantity being compared across structures is the marginal likelihood ---
the likelihood with the parameters integrated out --- and one cannot integrate a
coefficient out without a density to integrate it against. So each included
coefficient gets a prior, called the slab, and each excluded one is a point mass at
exactly zero. Conditional on $\sigma_i^2$ and on a common scale $g$:
\begin{eqnarray}
\alpha_i[\gamma_i] \;\bigm|\; \sigma_i^2, g &\sim&
  \mathcal{N}\bigl(0,\; g\,\sigma_i^2\, I_k\bigr),
  \Elabel{E-93A}{E93_A} \\
k &=& |\gamma_i|.
  \Elabel{E-93B}{E93_B}
\end{eqnarray}
The spike-and-slab pairing --- a point mass, not a narrow density --- is what
separates the output from ridge and from least squares, which put a small non-zero
value everywhere. Scaling by $\sigma_i^2$ makes $g$ dimensionless and keeps the
marginal in closed form.

\begin{recordbox}
\textbf{$p^M$ is not the whole penalty, and the part it is missing is the part that
does the work at large $T$.}

The structural term of \eqref{star} charges a \emph{fixed} amount for each included
entry --- $\log(1/p)$ nats, or after \S\ref{sec:rate} about $\log n_i$ nats. That
charge does not depend on the data and does not move with the sample size. The slab
supplies a second charge, and it is of a different kind: integrating an included
coefficient against \eqref{E93_A} spreads its prior mass over a range of values, and
the data then concentrates on a much narrower range, so the fraction of prior mass
that survives is small. By \S\ref{sec:marginal} that fraction is
$(1+g\mu_j)^{-1/2}$ per included direction, where $\mu_j$ grows linearly in $T$ ---
so the second charge is $\approx\frac12\log T$ \textbf{nats per included entry, and
it grows with the sample size}.

That growth is exactly what makes the method get \emph{more} selective as data
accumulates rather than less. Without it, the inclusion threshold on a squared
$t$-statistic would sit at a constant $2\log(1/p)$ forever, so a fixed fraction of
the null entries would clear it at every sample size and the false-positive count
would never fall. With it, the threshold rises like $\log T$ while the null
distribution of the statistic does not move, and the expected number of spurious
links falls like $1/\sqrt{T}$ --- which is \eqref{E133}, the quantity the whole
benchmark sweep is designed to test. A derivation carrying $p^M$ alone would look
right at any single sample size and be wrong across sample sizes.
\end{recordbox}

\begin{remark}[the entries are \emph{conditionally} i.i.d., and the distinction is
real]\label{rem:iid}
\eqref{E93_A} makes the included coefficients independent and identically distributed
\textbf{given $\sigma_i^2$ and} $g$, and the inclusion indicators independent and
identically distributed \textbf{given the rate} \eqref{E96_A}. Neither statement holds
marginally, and it would be careless to claim it did. The coefficients within one
equation share that equation's dispersion $\sigma_i^2$, and \emph{all} coefficients
in \emph{all} equations share the single slab scale $g$, which is inferred;
integrating $g$ out therefore couples every entry in the matrix to every other.
Likewise integrating the rate couples the indicators within an equation, which is
precisely how the multiplicity correction of \S\ref{sec:rate} arises --- the coupling
is the mechanism, not a defect. Nothing departs from conditional independence across
entries anywhere in the derivation: the only shared quantities are $\sigma_i^2$ (per
equation), $\rho_i$ (per equation) and $g$ (global), and each is named where it
enters. The practical consequence is stated at Remark \ref{rem:decouple}: $N$
independent per-equation searches, inside one shared one-dimensional quadrature, and
nothing else.
\end{remark}

\subsection{Why the slab scale cannot be fixed}\label{sec:slabfixed}

\begin{proposition}[Lindley--Bartlett: a fixed diffuse slab empties the
model]\label{prop:lb}
Let $\gamma$ and $\gamma\cup\{j\}$ differ by one entry, and hold $g$ fixed. Then,
with $\Zc_i$ as in \eqref{E101},
\begin{equation}
\lim_{g\to\infty}\frac{\Zc_i(\gamma\cup\{j\},g)}{\Zc_i(\gamma,g)} \;=\; 0,
\tag{E-94}\label{E94}
\end{equation}
the ratio decaying as $g^{-1/2}$.
So for a fixed slab taken wide enough, \emph{every} entry is switched off regardless
of the data, and the inferred sparsity is a function of $g$ alone.
\end{proposition}

\begin{proof}
From \eqref{E101},
$\Zc_i\propto|I_k+gS_\gamma|^{-1/2}(b_\sigma+\frac12R_\gamma(g))^{-(a_\sigma+T/2)}$.
As $g\to\infty$, $R_\gamma(g)$ tends to the ordinary residual sum of squares, finite
and independent of $g$, so the second factor tends to a finite positive limit. The
first satisfies $|I_{k+1}+gS_{\gamma\cup j}|/|I_k+gS_\gamma|\sim g\tilde\mu_j$ for
the residual eigenvalue $\tilde\mu_j>0$, so the ratio of first factors is
$\sim(g\tilde\mu_j)^{-1/2}$. The log Bayes factor therefore falls by $\frac12\log g$
without bound.
\end{proof}

A fixed \emph{narrow} slab is the mirror failure: every coefficient's prior mass
concentrates at zero, the inclusion Bayes factor approaches one, and the posterior
over structures approaches the prior. Either way the answer is a function of a
constant nobody derived. \textbf{Therefore $g$ is given a hyperprior and integrated
out} (D-09, D-24):
\begin{eqnarray}
\sqrt{g} &\sim& \text{Half-Cauchy}(0,1), \quad\text{equivalently}
  \Elabel{E-95A}{E95_A} \\
\pi(g) &=& \frac{1}{\pi\sqrt{g}\,(1+g)}.
  \Elabel{E-95B}{E95_B}
\end{eqnarray}
Three properties, each consumed below. It is \emph{proper}, so Assumption
\ref{ass:proper} holds and the integral of \S\ref{sec:quad} converges at the lower
end --- the reflex choice $\pi(g)\propto1/g$ is improper and diverges there. Its tail
is $g^{-3/2}$, heavy enough that the data and not the hyperprior decides the upper
end. Its median is $g=1$, which on standardized coordinates reads ``a slab of the
same scale as the innovation.'' The \emph{scale} of the half-Cauchy is nonetheless a
set constant and is listed in \S\ref{sec:constants}.

\subsection{Third refinement: marginalising the rate}\label{sec:rate}

The rate $p$ in \eqref{star} is written as though it were known. It is not, and
fixing it would be the single most damaging thing one could do to this method. A
fixed rate charges $\log\frac{1-p}{p}$ nats for each included entry --- a constant,
chosen by whoever chose $p$, and \textbf{independent of how many candidate entries
there were}. A problem with 3200 candidates would then be penalised at exactly the
same rate per entry as a problem with twelve, so the method would over-fit badly at
scale and would have a tuning knob sitting where its central claim ought to be.

Integrating the rate out repairs this, and the repair is not a patch but the source
of the method's most useful property. Give the rate a Beta prior and integrate; the
integral is a Beta function, and what emerges is a prior on the pattern that depends
on the count through $\binom{n_i}{k_i}^{-1}$ rather than through $p^{k_i}$. That is a
\textbf{multiplicity correction}: a pattern with $k_i$ non-zeros is charged the
logarithm of the number of ways $k_i$ entries could have been chosen out of the $n_i$
on offer, which is about $\log n_i$ nats per included entry. \textbf{The penalty
therefore scales with the number of candidate entries with nothing set by hand} ---
it tightens automatically as the problem grows, and it is what makes the
false-positive count in \eqref{E133} independent of the size of the matrix. Nobody
chose $\log n_i$; marginalising the rate produced it.

Concretely, each equation carries its own rate, all drawn from the same Beta:
\begin{eqnarray}
\gamma_i\mid\rho_i &\sim& \operatorname{Bernoulli}(\rho_i)^{\otimes n_i},
  \Elabel{E-96A}{E96_A} \\
\rho_i &\overset{\text{iid}}{\sim}& \operatorname{Beta}(a,b),
  \Elabel{E-96B}{E96_B} \\
a = b &=& 1,
  \Elabel{E-96C}{E96_C} \\
n_i &=& Np.
  \Elabel{E-96D}{E96_D}
\end{eqnarray}
Marginalising each $\rho_i$ gives the object that carries the whole of the complexity
control on the discrete side:

\begin{proposition}[the structure prior]\label{prop:sprior}
\begin{eqnarray}
\pi(\gamma_i) &=& \int_0^1 \rho_i^{k_i}(1-\rho_i)^{n_i-k_i}\,
  \frac{\rho_i^{a-1}(1-\rho_i)^{b-1}}{B(a,b)}\,d\rho_i \;=\;
  \frac{B(a+k_i,\; b+n_i-k_i)}{B(a,b)},
  \Elabel{E-111A}{E111_A} \\
k_i &:=& |\gamma_i\setminus\{0\}|,
  \Elabel{E-111B}{E111_B}
\end{eqnarray}
which depends on $\gamma_i$ only through the count $k_i$, and the joint structure
prior is the product
\begin{equation}
\pi(\gamma) \;=\; \prod_{i=1}^{N}\frac{B(a+k_i,\; b+n_i-k_i)}{B(a,b)}.
\tag{E-135}\label{E135}
\end{equation}
At $a=b=1$ the induced law on each count is uniform on $k_i\in\{0,\dots,n_i\}$.
\end{proposition}

\begin{proof}
The integral is the Beta function by definition. For the last claim,
$B(1+k,1+n-k)=k!(n-k)!/(n+1)!=\binom{n}{k}^{-1}/(n+1)$, so
$\pi(k)=\binom{n}{k}\pi(\gamma)=1/(n+1)$.
\end{proof}

\begin{corollary}[the multiplicity correction, exhibited]\label{cor:mult}
At $a=b=1$,
\begin{equation}
\log\pi(\gamma_i) \;=\; -\log(n_i+1) \;-\; \log\binom{n_i}{k_i}.
\tag{E-112}\label{E112}
\end{equation}
The first term is constant across structures and cancels. The second is the
multiplicity correction, and for $k_i\ll n_i$ it is $\approx k_i\log(n_i/k_i)+k_i$:
\textbf{roughly $\log n_i$ nats per included entry}, growing with the number of
candidates while the data does not.
\end{corollary}

\begin{corollary}[the price of one more entry]\label{cor:price}
\begin{equation}
\Delta_\pi(k_i) \;=\; \log\frac{\pi(\gamma_i\cup\{j\})}{\pi(\gamma_i)} \;=\;
\log\frac{a+k_i}{b+n_i-k_i-1} \;\xrightarrow[a=b=1]{\qquad}\;
\log\frac{k_i+1}{n_i-k_i} \;\approx\; -\log n_i.
\tag{E-113}\label{E113}
\end{equation}
$O(1)$ arithmetic in the count alone, which is what Proposition \ref{prop:flip}
consumes.
\end{corollary}

\noindent\textbf{Three readings, all of them consequences.}
\begin{enumerate}[label=(\alph*), topsep=0pt, itemsep=0pt]
\item \textbf{The penalty adapts to problem size with no tuning.} At $N=50$, $p=1$:
  $n_i=50$, so $\approx3.9$ nats per entry. At $p=2$: $n_i=100$, 4.6 nats. At $N=6$,
  $p=2$: $n_i=12$, 2.5 nats. Nobody set these; they are $\log n_i$, and they came
  from integrating the rate.
\item \textbf{Small data therefore concentrates mass on few non-zeros
  automatically.} An entry is included only if the data buys back $\log n_i$ nats
  through the likelihood, which \S\ref{sec:inclusion} turns into an explicit
  threshold.
\item \textbf{A fixed inclusion probability would not do this.} With $\rho$ fixed the
  charge is $\log\frac{\rho}{1-\rho}$, independent of $n_i$: the penalty would not
  tighten as the candidate set grew, and $\rho$ would be a tuning constant. That is
  D-08's argument, derived.
\end{enumerate}

\begin{recordbox}
\textbf{Self-flagged weakness (W-1, rank 1). Per-equation rate versus one global rate
is a substantive choice worth $\log N$ nats per entry --- 3.9 at $N=50$ --- and it
moves the predicted false-positive count by a factor of $N$.} Under the per-equation
rate ruled here the multiplicity correction charges $\log n_i=\log(Np)$; under a
single global rate over all $n=N^2p$ candidates it charges $\log(N^2p)$, and
\eqref{E133} falls by a factor $\approx N$. Table \ref{tab:b1} in \S\ref{sec:predict}
prints \emph{both} columns so that the size of the ruling is visible. The
per-equation form is what makes Remark \ref{rem:decouple} true and the search factor
cleanly; the global form is the more conservative. The derivation is written so that
substituting one for the other changes \eqref{E111_A}--\eqref{E113} and \eqref{E135}
and nothing else. \textbf{Confirmation requested.}
\end{recordbox}

\subsection{The prior over the lag order}\label{sec:lagprior}

Without a proper prior over $p$, the posterior over $p$ --- a headline output --- is
normalised over whatever set the search happened to visit, i.e.\ it is a function of
the search budget. So:
\begin{eqnarray}
\pi(p) &=& \frac{(1-\varpi)\,\varpi^{p}}{1-\varpi^{p_{\max}+1}},
  \Elabel{E-97A}{E97_A} \\
p &=& 0,1,\dots,p_{\max},
  \Elabel{E-97B}{E97_B} \\
\varpi &=& \tfrac12,
  \Elabel{E-97C}{E97_C}
\end{eqnarray}
a geometric law renormalised over the \emph{declared} range of Assumption
\ref{ass:cond}. It is not a cap on the model: $\pi(p)$ is proper on any range
including an unbounded one, and D-10 forbids a cap that changes the score, which a
proper prior is not. $\varpi=\frac12$ costs 0.69 nats per lag against the $N^2$ extra
parameters a lag brings, so it is deliberately weak --- its job is to make the sum
finite, not to choose $p$. \S\ref{sec:cannot}(iii) is the one case where that fails.

\subsection{Every constant that gets set, declared}\label{sec:constants}

\textbf{The previous draft claimed that no quantity in v1 requires a user-supplied
numerical value. That claim was false and is withdrawn.} The correct statement is
narrower and still worth making: \emph{no quantity that carries a scale is left to
the user} --- the innovation scale, the slab scale and the inclusion rate are all
integrated out rather than set --- but a handful of constants are set, and one of
them, an evidence threshold, sits underneath both the quadrature grid and the search
stopping rule and had never been named.

AUTOCLASS declares its own equivalent and gives a number.
\texttt{classes-c.text} 154--160: \emph{``Those with log marginals that differ by
less than 5 are are considered to be nearly equally probable''};
\texttt{search-c.text} 78--82 explains the reason, which is that its own computation
is not more accurate than that. We adopt the practice and state our own number, which
differs from theirs because our reason differs: our marginal is closed form up to one
smooth one-dimensional integral, so our arithmetic is far more accurate than
AUTOCLASS's, and what we are declaring is not a computational accuracy but a
\textbf{reporting resolution}.

\begin{definition}[the evidence threshold, and the chain hanging from
it]\label{def:thresh}
\begin{eqnarray}
\delta_{\text{report}} &=& 1\text{ nat}, \quad\text{hence}
  \Elabel{E-137A}{E137_A} \\
\epsilon_{\text{stop}} &\le& \tfrac{1}{10}\delta_{\text{report}},
  \Elabel{E-137B}{E137_B} \\
\Delta_{\text{trunc}} &\le& \tfrac{1}{10}\epsilon_{\text{stop}}.
  \Elabel{E-137C}{E137_C}
\end{eqnarray}
Log Bayes factors below $\delta_{\text{report}}$ are reported as ``not
distinguished'' and no claim is made from them. The search stops when no
single-entry move improves the score by more than $\epsilon_{\text{stop}}$; the
quadrature grid must bound its own truncation error by $\Delta_{\text{trunc}}$, by
\eqref{E108}, and must report the achieved value.
\end{definition}

\paragraph{Why one nat.} The smallest decision the model ever makes is the inclusion
of a single entry, and by \eqref{E113} that decision is fought over
$\log n_i\approx4$ nats at the benchmark size, $\approx8$ under a global rate. A
reporting threshold at one nat is a quarter of the smallest genuine decision, so it
cannot flip one; and as a bare odds statement $e:1$ is a preference nobody would call
evidence. The two factors of ten below it exist so that the stopping rule cannot
chase quadrature noise and the quadrature cannot chase arithmetic noise. \textbf{The
number is declared, not derived}, and a different declaration would change which
near-threshold entries are reported as decided --- not which are included.

The complete inventory. ``Effect'' is the magnitude by which the constant moves a
reported quantity, stated so that the weak ones can be recognised as weak without
being called derived.

\begin{center}
\small
\begin{tabular}{@{}>{\raggedright\arraybackslash}p{3.5cm}
                  >{\raggedright\arraybackslash}p{2.4cm}
                  >{\raggedright\arraybackslash}p{5.2cm}
                  >{\raggedright\arraybackslash}p{4.9cm}@{}}
\toprule
Constant & Where & Argument for the value & Effect if changed \\
\midrule
$a_\sigma=b_\sigma=\frac12$ & \eqref{E92_B} &
  ratio pinned by standardization, properness excludes 0; the level is set &
  $\approx0.05$ nats/entry at $T=200$; $\approx0.5$ at $T=20$ \\
half-Cauchy scale $=1$ & \eqref{E95_A} &
  median $g=1$ on standardized coordinates &
  small; $g$ is integrated, so this shifts the integrand not the answer \\
$a=b=1$ & \eqref{E96_C} &
  uniform on the rate; makes $\pi(k_i)$ uniform &
  $\approx0.06$ nats/entry \\
$\varpi=\frac12$ & \eqref{E97_C} & deliberately weak & 0.69 nats per lag \\
$\delta_{\text{report}}=1$ nat & \eqref{E137_A} &
  \textbf{declared}; a quarter of the smallest model decision &
  which near-threshold entries are called decided \\
$\epsilon_{\text{stop}}$, $\Delta_{\text{trunc}}$ & \eqref{E137_B}--\eqref{E137_C} &
  chained, one decade apart &
  below $\delta_{\text{report}}$ by construction \\
$p_{\max}$ & Assumption \ref{ass:cond} &
  \textbf{declared before fitting}; it sets the conditioning set &
  changes the effective sample, hence \emph{every} $\log\Zc$; Bayes factors at fixed
  $p_{\max}$ are unaffected \\
$\varepsilon$ & \eqref{E124_B} &
  a \textbf{reporting} convention for the advisory ceiling $p^*$ &
  changes advice only; enters no score \\
$n_{\text{restart}}$, $n_{\text{keep}}$ & \S\ref{sec:search} &
  search budget, open item O-3 &
  $n_{\text{keep}}$ changes the reported $\pi(k\mid y)$, $\pi(p\mid y)$ (W-4) \\
$d$ & D-05 & fixed a priori, not inferred &
  large; \S\ref{sec:cannot}; the one genuine modelling knob \\
\bottomrule
\end{tabular}
\end{center}

Two entries deserve emphasis because they were misrepresented before.
\textbf{$p_{\max}$ and $p^*$ are different objects.} $p_{\max}$ is declared, fixes
the conditioning set, and therefore enters every marginal likelihood; the previous
draft set it equal to the computed $p^*$, which made a reporting convention
$\varepsilon$ silently change the effective sample. That is repaired here by
declaring $p_{\max}$ and demoting $p^*$ to advice (\S\ref{sec:ceiling}). And $d$
\textbf{is a genuine knob with a large effect}, not a data description;
\S\ref{sec:cannot} is what it costs and open item O-5 is where it is decided.

\begin{remark}[the honest comparison with AUTOCLASS]\label{rem:autoclass}
The project's earlier claim that AUTOCLASS ``never asked the user for a scale'' is
\textbf{false}. \texttt{preparation-c.text} 357--376 makes a per-attribute
\texttt{error} a required property of every real attribute, describes it as
\emph{``a lower bound on the width of the normal distribution''}, and states that
\emph{``broad error estimates tend to limit the number of classes''}: one exogenous
number per real attribute, and it controls the answer to AUTOCLASS's headline
question. Everything else in its prior is either a hard-coded constant or read from
global database statistics, which its own authors call ``cheating'' in those words,
with a footnote explaining that a prior is supposed to be independent of the evidence
(TR \S2.2).

The honest framing is better than the false one. \textbf{AUTOCLASS takes one scale
per attribute and gets a complexity bound in return; ACV takes none and therefore has
no complexity bound of that kind} --- which is precisely why \S\ref{sec:degen}(D-f)
has to \emph{refuse} the constant-coordinate case that AUTOCLASS's \texttt{error}
absorbs. On the one axis where we are stricter, we are stricter because D-09 repairs
a limitation AUTOCLASS's authors named themselves.
\end{remark}


\section{The marginal likelihood, in closed form}\label{sec:marginal}

\textbf{[Requires: Gaussian coordinates, \eqref{E92_A}, \eqref{E93_A}, diagonal
$\Sigma$]} This section evaluates the second term of the score. It is the reduction
target of charter \S7 and the only part of v1 needing any numerical approximation at
all, and that approximation is one smooth integral in one variable.

\subsection{The per-equation marginal, given the slab scale}

Fix equation $i$ and structure $\gamma$, write $S_\gamma$, $s=s_i[\gamma]$, $c=c_i$,
and define
\begin{eqnarray}
\Lambda_\gamma(g) &:=& S_\gamma + \tfrac1g I_k \;\succ\; 0,
  \Elabel{E-99A}{E99_A} \\
\hat\alpha_\gamma(g) &:=& \Lambda_\gamma(g)^{-1}s,
  \Elabel{E-99B}{E99_B} \\
R_\gamma(g) &:=& c - s'\Lambda_\gamma(g)^{-1}s \;\ge\; 0.
  \Elabel{E-100}{E100}
\end{eqnarray}

\begin{proposition}[the closed form]\label{prop:closed}
\begin{equation}
\Zc_i(\gamma,g) \;=\; (2\pi)^{-T/2}\,\bigl|I_k + g\,S_\gamma\bigr|^{-1/2}\,
\frac{b_\sigma^{a_\sigma}}{\Gamma(a_\sigma)}\,
\Gamma\Bigl(a_\sigma+\tfrac{T}{2}\Bigr)\,
\Bigl(b_\sigma+\tfrac{R_\gamma(g)}{2}\Bigr)^{-(a_\sigma+T/2)}.
\tag{E-101}\label{E101}
\end{equation}
\end{proposition}

\begin{proof}
Write $\alpha$ for $\alpha_i[\gamma]$ and $\sigma^2$ for $\sigma_i^2$. The integrand
is
\begin{eqnarray}
e^{\ell_i}\,\pi(\alpha\mid\sigma^2,g)\,\pi(\sigma^2) &=& (2\pi\sigma^2)^{-T/2}
  \exp\bigl\{-\tfrac{1}{2\sigma^2}(c-2\alpha's+\alpha'S_\gamma\alpha)\bigr\}
  \nonumber\\
&&\times\,(2\pi g\sigma^2)^{-k/2}\exp\bigl\{-\tfrac{1}{2g\sigma^2}\alpha'\alpha\bigr\}
  \times\tfrac{b_\sigma^{a_\sigma}}{\Gamma(a_\sigma)}(\sigma^2)^{-a_\sigma-1}
  e^{-b_\sigma/\sigma^2}.
  \Elabel{E-98}{E98}
\end{eqnarray}
\emph{Complete the square in $\alpha$.} The two exponents combine to
$-\frac{1}{2\sigma^2}(c-2\alpha's+\alpha'\Lambda_\gamma\alpha)$, and since
$\Lambda_\gamma\succ0$,
$\alpha'\Lambda_\gamma\alpha-2\alpha's=(\alpha-\hat\alpha_\gamma)'\Lambda_\gamma
(\alpha-\hat\alpha_\gamma)-s'\Lambda_\gamma^{-1}s$, leaving $R_\gamma$ of
\eqref{E100} plus a centred quadratic. $R_\gamma\ge0$ because it equals
$\min_\alpha\{\sum_t(y_{ti}-\alpha'z_t)^2+g^{-1}\|\alpha\|^2\}$.

\emph{Integrate over $\alpha\in\mathbb{R}^k$.} The Gaussian integral gives
$(2\pi\sigma^2)^{k/2}|\Lambda_\gamma|^{-1/2}$; multiplying by the slab normaliser
$(2\pi g\sigma^2)^{-k/2}$ gives
$g^{-k/2}|\Lambda_\gamma|^{-1/2}=|g\Lambda_\gamma|^{-1/2}=|I_k+gS_\gamma|^{-1/2}$.
$\sigma^2$ \textbf{has cancelled from this factor entirely}, which is what scaling
the slab by $\sigma^2$ bought and is why the next step is clean.

\emph{Integrate over $\sigma^2$.} Substituting $w=(b_\sigma+R_\gamma/2)/\sigma^2$
turns the remaining integral into
$(b_\sigma+R_\gamma/2)^{-(a_\sigma+T/2)}\Gamma(a_\sigma+T/2)$, finite because
$a_\sigma+T/2>0$ and $b_\sigma>0$. Collecting gives \eqref{E101}.
\end{proof}

In log space, which is what is implemented:
\begin{multline}
\log\Zc_i(\gamma,g) = -\tfrac{T}{2}\log2\pi - \tfrac12\log\bigl|I_k+gS_\gamma\bigr|
 + a_\sigma\log b_\sigma - \log\Gamma(a_\sigma) \\
 + \log\Gamma\bigl(a_\sigma+\tfrac{T}{2}\bigr)
 - \bigl(a_\sigma+\tfrac{T}{2}\bigr)\log\bigl(b_\sigma+\tfrac{R_\gamma(g)}{2}\bigr).
\tag{E-102}\label{E102}
\end{multline}
\textbf{Only two terms depend on the structure}: the log-determinant, which is the
Occam factor in closed form, and $R_\gamma$, which is the fit. Everything else
cancels from every Bayes factor, so a structure comparison at fixed $g$ costs one
log-determinant and one quadratic form --- both delivered by
\eqref{E30_A}--\eqref{E86}.

\subsection[Reading (E-101)]{Reading \eqref{E101}: this \emph{is} the Occam factor}\label{sec:occam}

Let $\mu_1\ge\dots\ge\mu_k\ge0$ be the eigenvalues of $S_\gamma$. Then
$-\frac12\log|I_k+gS_\gamma|=-\frac12\sum_j\log(1+g\mu_j)$, so each included
regressor contributes $-\frac12\log(1+g\mu_j)$ nats: \textbf{zero} if it carries no
information ($\mu_j=0$), and $\approx-\frac12\log T$ once $g\mu_j\gg1$, since $\mu_j$
grows linearly in $T$. The volume correction therefore charges each \emph{effective}
parameter half a log of the sample size --- which is BIC's penalty, arrived at rather
than imposed, computed per direction rather than by counting, and automatically zero
for a direction the data does not identify. A parameter count cannot have that last
property. It is also why no separate ``effective dimension'' question arises in v1:
\eqref{E101} answers it direction by direction. Formally, with no hidden state there
is no symmetry group acting on the parameters, so the round-1 rank bound holds at
group dimension zero and
\begin{eqnarray}
D_{\eff} &=& D,
  \Elabel{E-21A}{E21_A} \\
\operatorname{rank} H_\ell(\hat\theta) &=& D \text{ generically},
  \Elabel{E-21B}{E21_B}
\end{eqnarray}
with $H_\ell:=-\nabla^2\ell$ the likelihood Hessian,
$\Lambda:=-\nabla^2\log\pi$ the prior curvature, and $H=H_\ell+\Lambda$ the posterior
Hessian, a distinction kept explicit at every use:
\begin{eqnarray}
H_\ell &=& -\nabla^2\ell,
  \Elabel{E-81A}{E81_A} \\
\Lambda &=& -\nabla^2\log\pi,
  \Elabel{E-81B}{E81_B} \\
H &=& H_\ell+\Lambda \;\succ\; 0,
  \Elabel{E-81C}{E81_C}
\end{eqnarray}
the positive definiteness because $H_\ell\succeq0$ (a Gram matrix over $\sigma_i^2$)
and $\Lambda=I/(g\sigma_i^2)\succ0$. \textbf{There is no singular-Hessian failure mode
in v1}, at any $T$, including $T<k$.

\begin{remark}[the general-prior statement, and why it is not used]\label{rem:laplace}
For a general prior the same content is Laplace's method,
\begin{equation}
\Zc(\Sc) \;=\; e^{\ell(\hat\theta)}\,\pi(\hat\theta\mid\Sc)\,(2\pi)^{D/2}\,
|H|^{-1/2}\,\bigl(1+O(T^{-1})\bigr),
\tag{E-18}\label{E18}
\end{equation}
\begin{equation}
\log\Zc(\Sc) \;=\; \underbrace{\ell(\hat\theta)}_{\text{fit}} \;+\;
\underbrace{\log\pi(\hat\theta\mid\Sc) + \tfrac{D}{2}\log2\pi
 - \tfrac12\log|H|}_{\log\Omega(\Sc),\text{ the Occam factor}} \;+\; O(T^{-1}),
\tag{E-19}\label{E19}
\end{equation}
whose second group of factors is the fraction of prior mass surviving contact with
the data --- the Occam factor of \eqref{star} written for a general prior. That is
what \eqref{E101} evaluates exactly. \textbf{Laplace is never invoked to compute a
reported number in v1}; it is retained here because it is the prior-agnostic
statement of what an Occam factor is, because the exponential-family interface of
\S\ref{sec:obs} needs it, and because dropping its bracket is where BIC comes from ---
and BIC would charge $\frac12\log T$ per entry and \emph{nothing} for the fact that
the entry was chosen out of $n_i$ candidates, discarding the multiplicity mechanism
and keeping a proxy.
\end{remark}

\subsection{The slab-scale quadrature}\label{sec:quad}

$g$ must be integrated, not maximised. Writing $u=\log g$,
\begin{equation}
\Zc(\Sc) \;=\; \int_0^\infty\Bigl[\prod_{i=1}^{N}\Zc_i(\gamma_i,g)\Bigr]\pi(g)\,dg
\;=\; \int_{-\infty}^{\infty} e^{\psi_\Sc(u)}\,du,
\tag{E-103}\label{E103}
\end{equation}
\begin{equation}
\psi_\Sc(u) \;=\; \sum_{i=1}^{N}\log\Zc_i(\gamma_i,e^u) \;+\; \log\pi(e^u) \;+\; u.
\tag{E-104}\label{E104}
\end{equation}
One dimension, the same integral for every structure. Differentiating \eqref{E102},
\begin{equation}
\psi_\Sc'(u) = \sum_{i=1}^{N}\Biggl[
\underbrace{-\tfrac12\sum_j\frac{g\mu_{ij}}{1+g\mu_{ij}}}_{\text{Occam, }\in[-k_i/2,0]}
+ \underbrace{\bigl(a_\sigma+\tfrac{T}{2}\bigr)
  \frac{\|\hat\alpha_{\gamma_i}(g)\|^2/(2g)}{b_\sigma+R_{\gamma_i}(g)/2}}_{\text{fit gain, }\ge0}
\Biggr] + \frac{d}{du}\bigl[\log\pi(e^u)+u\bigr],
\tag{E-105}\label{E105}
\end{equation}
which is the Occam mechanism visible in one derivative: the penalty saturates at
$k_i/2$ per equation while the gain is bought at a rate that falls as the slab
widens, and the mode is where they balance.

\begin{proposition}[the two ends]\label{prop:ends}
\begin{enumerate}[label=(\alph*), leftmargin=2em, topsep=0pt, itemsep=0pt]
\item $g\to0$. $\hat\alpha_\gamma\to0$, $R_\gamma\to c$, $|I+gS_\gamma|\to1$, so
\begin{equation}
\lim_{g\to0}\Zc_i(\gamma,g) \;=\; \Zc_i(\varnothing) \;=\;
(2\pi)^{-T/2}\tfrac{b_\sigma^{a_\sigma}\Gamma(a_\sigma+T/2)}{\Gamma(a_\sigma)}
\bigl(b_\sigma+\tfrac{c}{2}\bigr)^{-(a_\sigma+T/2)},
\tag{E-106}\label{E106}
\end{equation}
independently of $\gamma$: the slab collapses onto the spike and every structure
scores alike. The left tail is governed entirely by the hyperprior, so convergence
there is equivalent to $\pi$ being proper at 0 --- which is why $\pi(g)\propto1/g$ is
inadmissible.
\item $g\to\infty$. $R_\gamma\to c-s'S_\gamma^{+}s$ and
$|I+gS_\gamma|\sim g^{k^{\eff}}\prod_{\mu_j>0}\mu_j$, so
\begin{eqnarray}
\log\Zc_i(\gamma,g) &=& \mathit{const} - \tfrac{k_i^{\eff}}{2}u + O(e^{-u}),
  \Elabel{E-107A}{E107_A} \\
\psi_\Sc'(u) &\to& -\tfrac12\bigl(1+\textstyle\sum_i k_i^{\eff}\bigr),
  \Elabel{E-107B}{E107_B}
\end{eqnarray}
the Lindley--Bartlett slope of Proposition \ref{prop:lb}: the right tail decays
linearly in $u$, the more steeply the larger the structure.
\end{enumerate}
\end{proposition}

\begin{corollary}[a structure-independent floor]\label{cor:floor}
By (a),
$\Zc(\Sc)\ge[\prod_i\Zc_i(\varnothing)]\Pr(g\le\epsilon)(1-o(1))$ with a bound not
depending on $\Sc$. \textbf{No structure can score worse than the empty one by more
than $-\log\Pr(g\le\epsilon)$.} Two consequences for the grid: the floor cannot be
truncated away without changing Bayes factors, and it cannot be resolved by adding
nodes, because the integrand is flat there. Include it, sample it coarsely.
\end{corollary}

\paragraph{What makes a grid adequate.} Four conditions, in order of how easily each
is violated.
\begin{description}[leftmargin=2.5em, labelindent=0pt, topsep=0pt, itemsep=0pt,
  font=\normalfont]
\item[(Q1)] \textbf{The same grid for every structure compared}, or the score is a
  function of the grid and D-10 is violated by construction. This is the condition
  D-24's warning is about; the cache key \eqref{E109} enforces it.
\item[(Q2)] \textbf{Brackets every mode in the retained set}, not just the modal
  structure's. By \eqref{E105} the mode moves with $\sum_ik_i^{\eff}$ and with the
  signal; the bracket follows from setting $\psi'=0$ at the extreme structures.
\item[(Q3)] \textbf{Resolves the curvature.} The Occam term's second derivative is
  at most $\frac18\sum_ik_i$ in magnitude, so the peak half-width in $u$ is
  $\Omega((\sum_ik_i)^{-1/2})$ --- narrower for larger structures, but only as a
  square root. Gauss--Legendre with $n_Q$ nodes on a bracket of width $W$ resolves
  features of scale $W/n_Q$, so a fixed wide bracket at $N=50$, $k_i\approx7$,
  $W\approx20$ needs $n_Q\approx370$, and a bracket centred adaptively on the mode
  needs far fewer. The implementation must report which it did and the achieved
  accuracy; the previous draft quoted $n_Q=32$ in a cost claim and $\gtrsim370$ here
  without reconciling them.
\item[(Q4)] \textbf{Truncation error bounded, in nats, and reported.}
\begin{equation}
\Delta_{\text{trunc}} \;\le\; \Pr\bigl(g\le e^{u_{\min}}\bigr) \;+\;
\frac{\exp\{\psi_\Sc(u_{\max})\}}{\frac12\bigl(1+\sum_ik_i^{\eff}\bigr)},
\tag{E-108}\label{E108}
\end{equation}
  the first term from Corollary \ref{cor:floor} and the second by integrating the
  linear tail of \eqref{E107_A}. \textbf{This must be at most $\Delta_{\text{trunc}}$
  of} \eqref{E137_C}, and doubling $n_Q$ must change every reported $\log\Zc$ by less
  than that. A grid chosen until the results looked right is a tuning constant by
  another route.
\end{description}

\paragraph{The cache.} The nodes are fixed once. Each $\log\Zc_i(\gamma_i,e^{u_q})$
depends only on $(\gamma_i,c_i,s_i,S_{zz},T)$, so the key is
\begin{equation}
\text{key} \;=\; (i,\;\gamma_i,\;q),
\tag{E-109}\label{E109}
\end{equation}
and a move that changes $\gamma_i$ alone invalidates $n_Q$ entries for one $i$ and
none for the other $N-1$.


\section{When does an entry get switched on?}\label{sec:inclusion}

Everything so far assembles into one number per candidate move. This section
computes it.

The model-comparison machinery is three lines, and every other section either
evaluates $\Zc$ or specifies $\pi$:
\begin{equation}
\pi(y,\theta\mid\Sc) \;=\; f(y\mid\theta,\Sc)\,\pi(\theta\mid\Sc),
\tag{E-14}\label{E14}
\end{equation}
\begin{equation}
\Zc(\Sc) \;=\; \int_{\Theta_\Sc} f(y\mid\theta,\Sc)\,\pi(\theta\mid\Sc)\,d\theta,
\tag{E-15}\label{E15}
\end{equation}
\begin{equation}
\pi(\Sc\mid y) \;=\; \frac{\pi(\Sc)\,\Zc(\Sc)}{\sum_{\Sc'}\pi(\Sc')\,\Zc(\Sc')},
\tag{E-16}\label{E16}
\end{equation}
with $y$ everywhere abbreviating $y_{1:T}$ conditional on $y_{1-p_{\max}:0}$
(Assumption \ref{ass:cond}). Both \eqref{E15} and \eqref{E16} require Assumption
\ref{ass:proper}; \eqref{E16} additionally requires the candidate set to be a set,
which is what \eqref{E97_A} supplies. Writing
\begin{eqnarray}
\operatorname{score}(\Sc) &=& \log\pi(\Sc)+\log\Zc(\Sc),
  \Elabel{E-24A}{E24_A} \\
\pi(\Sc\mid y) &\propto& e^{\operatorname{score}(\Sc)}.
  \Elabel{E-24B}{E24_B}
\end{eqnarray}
\textbf{Both terms carry complexity control and they do different jobs}:
$\log\pi(\Sc)$ supplies the \emph{multiplicity} correction --- how many structures of
a given size exist --- and the Occam factor inside $\log\Zc$ supplies the
\emph{volume} correction --- how much prior mass a structure's continuous parameters
waste. A derivation carrying one and not the other is wrong in a way invisible at
fixed $k$ and fatal across $k$, which is the only comparison this project makes.
\S\ref{sec:degen}(D-d) exhibits a case where the two disagree, which is the cleanest
demonstration that both are needed.

Let $\tilde\mu_j:=S_{jj}-S_{j\gamma}\Lambda_\gamma^{-1}S_{\gamma j}$ be the residual
information in candidate $j$ after the included regressors, $\hat\beta_j$ its partial
regression coefficient, and
\begin{equation}
z_j^2 \;:=\; \frac{\tilde\mu_j\,\hat\beta_j^2}{\sigma_i^2}
\tag{E-118}\label{E118}
\end{equation}
the squared partial $t$-statistic.

\begin{proposition}[the inclusion rule]\label{prop:incl}
The change in \eqref{E24_A} on adding entry $j$ to equation $i$ is
\begin{equation}
\Delta_j \;=\; \underbrace{-\tfrac12\log\bigl(1+g\tilde\mu_j\bigr)}_{\textit{Occam (volume)}}
\;+\; \underbrace{\tfrac12\,\frac{g\tilde\mu_j}{1+g\tilde\mu_j}\,z_j^2}_{\textit{fit}}
\;+\; \underbrace{\Delta_\pi(k_i)}_{\textit{multiplicity}, \text{ \eqref{E113}}},
\tag{E-119}\label{E119}
\end{equation}
so in the regime $g\tilde\mu_j\gg1$ the entry is included exactly when
\begin{equation}
\tcboxmath{ \begin{array}[t]{c}
z_j^2 \;>\; \tau \;:=\; \log\bigl(g\tilde\mu_j\bigr) + 2\log\dfrac{n_i-k_i}{k_i+1}
\;\approx\; \log T + 2\log\dfrac{n_i-k_i}{k_i+1} + O(1).
\end{array}}
\tag{E-120}\label{E120}
\end{equation}
\end{proposition}

\begin{proof}
The determinant term of \eqref{E102} changes by $-\frac12\log(1+g\tilde\mu_j)$,
since $\tilde\mu_j$ is exactly the Schur complement in the bordered determinant. The
residual term changes by \eqref{E86}: $R$ falls by
$\tilde\mu_j\hat\beta_j^2\cdot g\tilde\mu_j/(1+g\tilde\mu_j)$, which divided by
$2\sigma_i^2$ and multiplied by the prefactor
$(a_\sigma+T/2)/(b_\sigma+R/2)\to1/\sigma_i^2$ gives the fit term. The prior term is
\eqref{E113}. Setting $\Delta_j>0$ and using $g\tilde\mu_j/(1+g\tilde\mu_j)\to1$
gives $z_j^2>\log(1+g\tilde\mu_j)-2\Delta_\pi(k_i)$, and $\tilde\mu_j=\Theta(T)$.
\end{proof}

\begin{recordbox}
\textbf{Correction, objection O-004.} The previous draft wrote the second term of
\eqref{E120} as $2\log n_i$ --- that is, evaluated it at $k_i=0$ --- and computed
every downstream number from that form, while the structure it was describing carries
$k_i\approx7$. The approximation $k\ll n$ licenses dropping $k$ from $n-k$; it does
\textbf{not} license dropping $k+1$ from the denominator. At $n_i=50$, $k_i=7$ the
difference is $2\log8=4.16$ nats, and $\tau$ is exponentiated in a tail probability,
so the error is a factor of $e^{2.08}\approx8$ in every false-positive count. Every
dependent number in \S\ref{sec:predict} has been recomputed.
\end{recordbox}

\noindent Two readings of \eqref{E120}, both consequences.
\begin{enumerate}[label=(\alph*), topsep=0pt, itemsep=0pt]
\item \textbf{The threshold rises with the number of candidates and with the sample
  size, and falls as the structure fills up.} The $\log T$ is the volume correction;
  the $2\log((n_i-k_i)/(k_i+1))$ is the multiplicity correction. The second is the
  one that makes the method scale: doubling $N$ at fixed $T$ raises the bar
  automatically, with nothing set by hand.
\item \textbf{It contains no reference to \emph{which} entry $j$ is.} The threshold
  for the entry linking coordinate 3 to coordinate 47 is the threshold for the entry
  linking coordinate 3 to coordinate 4. That is the boxed property of page one,
  appearing where it matters.
\end{enumerate}


\docpart{Part B --- The worked example}

\section{Benchmark B1}\label{sec:bench}

Benchmark B1 (\texttt{workflow/benchmark.md}) is a VAR(1) and belongs here unchanged.
It is deliberately non-economic and small enough to print in full. \textbf{This
section derives what the machinery predicts before the run; agent 3 runs it.}

\subsection{The generator}\label{sec:gen}

$N=50$ coordinates arranged in ten groups of five; $p=1$;
$A_1\in\mathbb{R}^{50\times50}$ with $n=2500$ candidate entries. Ten dense $5\times5$
diagonal blocks (250 entries) with block spectral radius 0.7; a coupling budget of
$C=100$ entries confined to ten off-diagonal block positions on a directed cycle
$1\to2\to\dots\to10\to1$; $\rho(A_1)$ rescaled after assembly. Total 350 non-zeros
of 2500 --- 14\% density, about seven per equation. $\Sigma$ diagonal; $x_0$ from the
stationary distribution. The generator produces the \emph{differences} as the VAR and
cumulates to levels, so that the pipeline's $d=1$ is correct pre-processing --- which
by Corollary \ref{cor:overdiff} is not a nicety but the condition under which B1 is
inside the model class at all.

\begin{recordbox}
\textbf{What B1 is and is not for.} The generator's truth happens to be clustered.
\textbf{The model cannot see that and is not being asked to exploit it.} B1-C
(clustered) and B1-S (the same coupling budget scattered at random) are co-equal
headline cases, and the required result is that the method recovers whatever is there
in \emph{both} --- approximately block-structured from clustered data, scattered from
scattered data. A method that reported block structure on B1-S would have an
assumption baked in; that failure mode is what B1-S exists to expose, and with the
partition model deleted there is no longer any mechanism by which it could occur.
B1-C is therefore an illustration that parsimony may land on something block-like,
\emph{not} a target to reproduce.
\end{recordbox}

\subsection{What the machinery predicts}\label{sec:predict}

By \eqref{E121_A} below a true entry with population partial signal-to-noise ratio
$\eta_j$ (so that $\E[z_j^2]\approx T\eta_j$) is detected once
\begin{eqnarray}
T\eta_j &\gtrsim& \tau, \quad\text{equivalently}
  \Elabel{E-121A}{E121_A} \\
T &\gtrsim& \tau/\eta_j,
  \Elabel{E-121B}{E121_B} \\
\tau &=& \log T + 2\log\frac{n_i-\bar k}{\bar k+1},
  \Elabel{E-121C}{E121_C}
\end{eqnarray}
and for a \emph{null} entry $z_j^2\sim\chi^2_1$, so the expected count of spurious
links per equation at threshold $\tau$ is
\begin{equation}
\E[\text{FP per equation}] \;\approx\; n_i\Pr(\chi^2_1>\tau) \;\approx\;
n_i\sqrt{\tfrac{2}{\pi\tau}}\;e^{-\tau/2} \;=\;
n_i\sqrt{\tfrac{2}{\pi\tau}}\;\frac{\bar k+1}{(n_i-\bar k)\sqrt{T}}.
\tag{E-122}\label{E122}
\end{equation}

\begin{corollary}[the headline prediction, corrected]\label{cor:headline}
Summing \eqref{E122} over the $N$ equations and using $n_i/(n_i-\bar k)\to1$,
\begin{equation}
\tcboxmath{ \begin{array}[t]{c}
\E[\text{FP total}] \;\approx\; N\,(\bar k+1)\,\sqrt{\dfrac{2}{\pi\tau}}\;
\dfrac{1}{\sqrt{T}}
\end{array}}
\tag{E-133}\label{E133}
\end{equation}
--- \textbf{in which the number of candidate entries $n_i$ has cancelled}: the
multiplicity correction raises the threshold by precisely enough to absorb it. The
expected number of spurious links does not depend on how many entries were on offer.
\end{corollary}

\subsection[Not parameter-free]{This prediction is not parameter-free, and it was described as if it
were}\label{sec:notfree}

\begin{recordbox}
\textbf{The cancellation of $n_i$ is real and it is the cleanest available statement
of what this method buys. The claim that \eqref{E133} is ``a prediction with no free
parameter in it'' is false and is withdrawn.} What survives the cancellation is
$N(\bar k+1)$, and $\bar k$ --- the mean number of entries retained per equation ---
is an \emph{output of the fit}, not a property of the design. \eqref{E133} is
therefore a prediction \emph{conditional on the reported sparsity}: it says that once
the method has told you it retained $\bar k$ entries per equation, you may predict
how many of them are spurious, and you may check that prediction against the sweep.
That is still falsifiable and still useful, but it is a self-consistency statement
rather than an a-priori one, and the difference matters because the supervisor told
Ben it was parameter-free and that was wrong.
\end{recordbox}

Table \ref{tab:b1} evaluates \eqref{E121_A}--\eqref{E133} at $\bar k=7$, $n_i=50$,
$N=50$. The last column prints the same quantity under a single \emph{global}
inclusion rate over all $n=N^2p=2500$ candidates, so that the size of the W-1 ruling
is visible: it is a factor of about $N$.

\begin{table}[ht]
\caption{What B1 predicts. $\tau$ from \eqref{E121_C}; $\eta$ needed is
$\tau/(\text{usable})$; FP columns from \eqref{E133}. \textbf{All values carry
$\bar k=7$.}}\label{tab:b1}
\centering
\small
\begin{tabular}{rrrrrrl}
\toprule
$T$ & usable & $\tau$ & $\eta_j$ needed & $\E$[FP]/eq. & $\E$[FP] total &
  (global rate) \\
\midrule
20   & 18   & 6.36  & 0.353   & 0.658  & 32.9 & 0.376 \\
50   & 48   & 7.28  & 0.152   & 0.389  & 19.5 & 0.231 \\
100  & 98   & 7.97  & 0.0813  & 0.263  & 13.1 & 0.160 \\
200  & 198  & 8.66  & 0.0437  & 0.178  & \textbf{8.92} & \textbf{0.110} \\
500  & 498  & 9.58  & 0.0192  & 0.107  & 5.36 & 0.068 \\
1000 & 998  & 10.27 & 0.0103  & 0.0732 & 3.66 & 0.047 \\
2000 & 1998 & 10.96 & 0.00549 & 0.0501 & 2.51 & 0.033 \\
5000 & 4998 & 11.88 & 0.00238 & 0.0305 & 1.52 & 0.020 \\
\bottomrule
\end{tabular}
\end{table}

\paragraph{Read the two ends against each other.} At $T=20$ an entry needs
$\eta_j\gtrsim0.35$, which almost no entry of a 14\%-dense matrix with $\rho(A)=0.7$
will have: \emph{predicted, a nearly empty matrix}, which is the correct answer at
that sample size. The failure to watch for is a confident report of many links. At
$T=5000$ an entry needs only $\eta_j\gtrsim0.0024$, so essentially every true entry
is found while \eqref{E133} predicts about 1.5 spurious links across the whole
$50\times50$ matrix: \emph{predicted, the count plateaus at $\approx350$ and does not
climb toward 2500.} \textbf{A method that recovers the truth at $T=200$ and reports
550 links at $T=5000$ has a broken multiplicity correction, and \eqref{E133} is the
number to check it against.} The $T^{-1/2}$ decay is itself the test: the threshold
rises like $\log T$ while the null distribution does not move.

\paragraph{Two caveats on the table, stated rather than buried.} At $T=20$ the usable
sample is 18 against $n_i=50$ candidates, so most structures in the space have
$k_i>T$; \eqref{E128} keeps the marginal likelihood finite there, but the $\chi^2_1$
approximation behind \eqref{E122} is not reliable and the entry `32.9' should be read
as ``tens, not units'' rather than as a number. And the coupling entries' magnitudes
are not pinned by \texttt{benchmark.md}: since $\eta_j\propto a_{ij}^2$, the $T$ at
which the inter-group coupling becomes detectable is a free choice of whoever writes
the generator. At a coupling magnitude of 0.1 the implied detection $T$ is of order
$10^3$, several times the charter's target. \textbf{\texttt{benchmark.md} needs a
magnitude rule before the sweep is run}, and item (G) below is how it gets checked.

\subsection[Format for the generating parameters]{The appendix format for the generating parameters}\label{sec:appfmt}

Fixed here so that agent 3 and agent 5 agree in advance; produced \emph{by running
the generator}, so the printed numbers and the numbers used are the same by
construction. (A) ten within-group blocks, one $5\times5$ table each, six significant
figures, each with its own $\rho$ before global rescaling. (B) coupling entries as
$(i,j,\text{value})$ triples grouped by block-edge, with a count per edge and a total
equal to $C=100$. (C) $\rho(A_1)$ before and after rescaling, and the factor. (D)
$\Sigma$, $\sigma_0$, and $\sigma_0/(\text{stationary s.d.})$. (E) RNG seed and
generator version, verbatim. (F) the $10\times10$ block-level occupancy matrix with
the strong-connectivity check and its result. \textbf{(G) the derived quantities this
document predicts} --- per-entry $\eta_j$ as a histogram \emph{split by within-block
and coupling}, the implied detection $T$ per entry, and the predicted $\E$[FP] at
each $T$ in the sweep.

Item (G) is this document's one addition to \texttt{benchmark.md} \S6, and it is what
lets the run \emph{refute} \S\ref{sec:predict} rather than merely report numbers.


\docpart{Part C --- General lag order, and the rest of the machinery}

\section{General lag order}\label{sec:lag}

\subsection{What changes, and what does not}

Exactly three things change, and all three are dimension bookkeeping. The regressor
vector lengthens to $z_t\in\mathbb{R}^{1+Np}$ by \eqref{E134_A}, so $S_{zz}$ is
$(1+Np)\times(1+Np)$ and $n_i=Np$. The parameter count grows to $D=N^2p+2N$, so each
additional lag costs $N^2$ parameters --- \textbf{not one}, which is the intuition a
reader arrives with from univariate Box--Jenkins. And the conditioning set lengthens,
with $p_{\max}$ fixed before fitting (Assumption \ref{ass:cond}).

Everything else is unchanged, for a structural reason rather than a list of
coincidences: \textbf{from \eqref{E134_A} onward, $p$ enters only through the length of
$z_t$.}

\begin{proposition}[the $p=1$ results generalise by substitution]\label{prop:subst}
With $z_t\in\mathbb{R}^{1+Np}$: \eqref{E26_A}--\eqref{E27} and Proposition
\ref{prop:suff} hold verbatim, sufficiency being a statement about the regression
form and indifferent to how $z_t$ was assembled; \eqref{E28_A}--\eqref{E30_C} and
Propositions \ref{prop:flip} hold verbatim, the decoupling using only diagonality of
$\Sigma$ and the bordering using only $\Lambda_\gamma\succ0$; Proposition
\ref{prop:closed} and \eqref{E101}--\eqref{E102} hold verbatim, their proof never
using $\dim z_t$; \S\ref{sec:quad} holds verbatim, $p$ entering only through
$k_i^{\eff}$; and \S\ref{sec:rate} holds with $n_i=Np$, so \textbf{the multiplicity
correction tightens with $p$ automatically, by $\log p$ nats per entry}.
\end{proposition}

\begin{proof}
Each clause is the observation that the cited proof refers to $z_t$ only through its
dimension and to $\bbA$ only through its row count. The last needs one line:
$\Delta_\pi\approx-\log n_i=-\log(Np)=-\log N-\log p$.
\end{proof}

The one place $p$ does more than lengthen a vector is collinearity: the columns of
$z_t$ at different lags are lagged copies of one another and are correlated by
construction, so the near-collinear regime is \emph{more} likely at $p=2$ than at
$p=1$. That is why \S\ref{sec:cond} is written as a diagnostic to be run rather than a
caveat to be noted.

\paragraph{Comparability.} \emph{Across $p$: comparable.} Every candidate is a
density for the same $y_{1:T}$ given the same conditioning set; a candidate with
$p<p_{\max}$ simply does not use the earlier pre-sample values, and $\pi(p)$ of
\eqref{E97_A} makes the posterior proper. \emph{Across $d$: not comparable} (D-05):
different differencing orders are models for different data, and the missing
Jacobian and effective-sample correction are not supplied here. \S\ref{sec:cannot}
is where that bites.

\subsection[The lag ceiling is advice]{The reported lag ceiling is advice, not a cap}\label{sec:ceiling}

D-23 requires that the practical limit on $p$ be shown to be a property of the data
volume and not of the theory. The program therefore computes and prints
\begin{eqnarray}
p^*(N,T;\varepsilon) &=& \max\Bigl\{\,p\in\mathbb{N} \;:\;
  \frac{N\,p}{T-p} \;\le\; \bigl(1-\sqrt{\varepsilon}\bigr)^2\Bigr\},
  \Elabel{E-124A}{E124_A} \\
\varepsilon &=& \tfrac14\text{ by default},
  \Elabel{E-124B}{E124_B}
\end{eqnarray}
from a Marchenko--Pastur aspect-ratio condition: at full inclusion the smallest
eigenvalue of $S_\gamma$ is $\approx(T-p)(1-\sqrt{c})^2$ with $c=n_i/(T-p)$, and
requiring the smallest direction to retain a fraction $\varepsilon$ of the average
information gives $c\le(1-\sqrt{\varepsilon})^2$. Alongside it the program prints the
search-space size
\begin{equation}
\log_2\bigl|\{\text{structures}\}\bigr| \;=\; N^2p\text{ bits},
\tag{E-125}\label{E125}
\end{equation}
so the user can see what the restart budget is being asked to cover.

\begin{recordbox}
\textbf{Correction, objection O-003.} The previous draft set $p_{\max}=p^*$. Since
$p_{\max}$ fixes the conditioning set and the effective sample, that made the
reporting convention $\varepsilon$ change \emph{every} marginal likelihood in the
document, while the text asserted it changed none. Here $p_{\max}$ is
\textbf{declared} (\S\ref{sec:constants}) and $p^*$ is \textbf{advisory}: it enters
no score and no normaliser. It tells the user where the posterior over structures is
becoming too diffuse to read, which is a different statement from where the model
breaks.
\end{recordbox}

Worked values, so a reader with more data can compute their own: $N=40,T=200$ gives
$p^*=1$ at $\varepsilon=\frac14$ and 2 at $\varepsilon=\frac1{16}$; $N=50,T=200$
gives 0 and 2; $N=50,T=5000$ gives 24 and 55; $N=6,T=20000$ gives 800 and 1714; and
$N=1,T=200$ gives 40 and 72.

Three observations. The last row shows the criterion \emph{recovering univariate
practice}: at $N=1$, \eqref{E124_A} reduces to $p\le T/5$, which is why Box--Jenkins
runs $p=3$ or 4 without a second thought and why that intuition does not survive
multiplication by $N$. Second, \textbf{at $\varepsilon=\frac14$ neither the charter's
target configuration nor the benchmark's clears the range it is specified at}, and I
report the disagreement rather than back-solve for the $\varepsilon$ that would
remove it. Third, the criterion is deliberately pessimistic because the search never
forms the full $Np\times Np$ system --- it works at $|\gamma_i|\approx7$, where the
aspect ratio is $7/198$ and the design is comfortable at any $\varepsilon$. The
resolution recommended is to report both: \eqref{E124_A} as the pre-fit advisory range
and the post-fit aspect ratio $\bar k/(T-p)$ alongside the structure. Both are
reports, so D-10 is untouched either way.

\subsection{Instantiation at $N=2$, $p=2$}\label{sec:inst}

\textbf{Nothing has ever been checked at $p=2$ anywhere in this project}, and the
whole operating range is $p\in\{1,2\}$. At $N=2$, $p=2$, $T=3$, with symbolic
$y_{-1},y_0,y_1,y_2,y_3\in\mathbb{R}^2$ and $z_t=(1,y_{t-1}',y_{t-2}')'$,
\begin{eqnarray}
S_{zz} &=& \sum_{t=1}^{3} z_t z_t',
  \Elabel{E-123A}{E123_A} \\
S_{xz} &=& \sum_{t=1}^{3} y_t z_t',
  \Elabel{E-123B}{E123_B} \\
S_{xx} &=& \sum_{t=1}^{3} y_t y_t'.
  \Elabel{E-123C}{E123_C}
\end{eqnarray}
Note the structure a $p=1$ check cannot exhibit: $S_{zz}$ has \emph{repeated blocks},
because the same observations appear at two lags. That is the algebraic source of lag
collinearity and it is invisible at $p=1$. The checks, in priority order:
\eqref{E27} with a \textbf{full symmetric} $\Sigma$, then \eqref{E28_A} under imposed
diagonality; \eqref{E101} at $\gamma_i$ containing one regressor from each lag,
against the direct double integral \eqref{E98}; \eqref{E30_A} and \eqref{E86} for a move
across a lag boundary; and \eqref{E119}'s Schur complement when $j$ is a lag-2
regressor and $\gamma$ already contains its lag-1 counterpart --- the place lag
collinearity first shows in the score, and it has never been computed.


\section{Search}\label{sec:search}

The score, instantiated, is
\begin{equation}
\operatorname{score}(\Sc) = \underbrace{\log\pi(p)}_{\text{\eqref{E97_A}}}
+ \underbrace{\sum_{i=1}^{N}\log\frac{B(a+k_i,\;b+n_i-k_i)}{B(a,b)}}_{\text{\eqref{E135}}}
+ \underbrace{\log\int_{-\infty}^{\infty}e^{\psi_\Sc(u)}\,du}_{\text{\eqref{E103}, fixed grid}}.
\tag{E-126}\label{E126}
\end{equation}
\textbf{No other term exists.} A reviewer auditing for D-10 violations checks that
\eqref{E126} is the only quantity ever compared, anywhere.

\paragraph{Moves.} Add $j$ to $\gamma_i$; delete $j$ from $\gamma_i$; swap
$j\to j'$ within $\gamma_i$. All cost $O(k_i^2)$ per quadrature node by Proposition
\ref{prop:flip}. Index 0, the intercept, is never a move.

\paragraph{What is held.} For each retained $(i,q)$: the Cholesky factor of
$S_{\gamma_i}+e^{-u_q}I$, updated by \eqref{E30_A}; the running scalar
$s_i[\gamma_i]'\Lambda^{-1}s_i[\gamma_i]$, updated by \eqref{E86}; and the counts
$k_i$, consumed by \eqref{E113} in $O(1)$. The outer integral is then re-evaluated by
summing $n_Q$ cached per-equation values, or $O(n_Q)$ if the sum
$\sum_i\log\Zc_i(\gamma_i,e^{u_q})$ is itself kept running. \textbf{That is the
entire cost of the one remaining coupling between equations: one running sum per
node.}

\paragraph{Stopping and reporting.} Sweep all single-flip neighbours of every
equation, accept the best improving move, stop when a full sweep yields no move worth
more than $\epsilon_{\text{stop}}$ nats \eqref{E137_B}. Restart $n_{\text{restart}}$
times from structures drawn from the prior \eqref{E135}, seeded deterministically.
Retain the best $n_{\text{keep}}$ distinct structures over all restarts and all
$p\le p_{\max}$; reported quantities are computed by normalising \eqref{E16} over the
retained set.

\begin{recordbox}
\textbf{Self-flagged weakness (W-4, rank 4).} Normalising over the retained set
approximates normalising over all structures and the error is not bounded here. A
proper $\pi(p)$ makes the \emph{idealised} posterior exist; it does not make the
\emph{computed} one budget-independent. The cheap honest mitigation is to report the
total retained mass against the same quantity at half the restart budget: if the two
differ materially the search has not converged and the posterior should not be
quoted. That is a report, not a score.
\end{recordbox}


\section{Degenerate cases}\label{sec:degen}

\begin{description}[style=nextline, leftmargin=2em, topsep=0pt]
\item[(D-a) An all-zero row, $\gamma_i=\{0\}$.]
  Coordinate $i$ has no incoming links. Then $k_i=0$, and \eqref{E101} gives
\begin{equation}
\Zc_i(\{0\}) = (2\pi)^{-T/2}\tfrac{b_\sigma^{a_\sigma}\Gamma(a_\sigma+T/2)}{\Gamma(a_\sigma)}
\bigl(b_\sigma+\tfrac{R_{\{0\}}}{2}\bigr)^{-(a_\sigma+T/2)}.
\tag{E-127}\label{E127}
\end{equation}
  \textbf{This is a perfectly well-defined model} --- coordinate $i$ is i.i.d.\ noise
  about its own drift --- and it is the correct answer at $T=20$. No special case is
  needed in the code: the empty structure is a structure. It is also the $p=0$ model
  of \S\ref{sec:red2}.

\item[(D-b) An all-zero column.]
  Coordinate $j$ appears in no equation's $\gamma_i$: it is not Granger-causal for
  anything at the retained structure. This is \emph{not} a degeneracy of any single
  equation --- every $\Zc_i$ is regular --- but it does mean row sparsity and column
  sparsity are different summaries and both must be printed, since an analyst reading
  only row counts will not see it. At $p\ge2$, ``no outgoing links'' must mean all
  $p$ of them.

\item[(D-c) $T$ small relative to $k$.]
  If $k_i\ge T$ then $S_\gamma$ is singular and the \emph{likelihood} is unbounded as
  $\sigma_i^2\to0$. At $N=40$, $p=2$, $T=20$ that is most structures in the space. It
  costs nothing:

  \begin{proposition}[the marginal likelihood is finite regardless]\label{prop:finite}
  For every $\gamma$, every $g>0$ and every $T\ge1$,
  \begin{equation}
  \Zc_i(\gamma,g) \;\le\; (2\pi)^{-T/2}\,
  \frac{b_\sigma^{a_\sigma}\Gamma(a_\sigma+T/2)}{\Gamma(a_\sigma)}\,
  b_\sigma^{-(a_\sigma+T/2)},
  \tag{E-128}\label{E128}
  \end{equation}
  uniformly over $\gamma$ and $g$.
  \end{proposition}

  \begin{proof}
  $\Lambda_\gamma\succ0$ for every $g>0$ regardless of $\operatorname{rank}S_\gamma$,
  so $R_\gamma\ge0$; $|I_k+gS_\gamma|\ge1$ since $gS_\gamma\succeq0$; and
  $b_\sigma+R_\gamma/2\ge b_\sigma>0$. Substituting into \eqref{E101} gives
  \eqref{E128}.
  \end{proof}

  The objection this answers predicted that at $T=20$ the score would be a function
  of the hyperparameter $b_\sigma$ rather than of the data. \textbf{That prediction
  assumed Laplace.} v1 computes \eqref{E101} exactly and never evaluates the
  likelihood at a maximiser, so what enters the score is $R_\gamma$, an integrated
  quantity bounded between 0 and $c_i$.

\item[(D-d) Duplicate or collinear regressors.]
  Suppose $z_{tj}=z_{tj'}$ for all $t$. Then
  $\Zc_i(\gamma\cup\{j\},g)=\Zc_i(\gamma\cup\{j'\},g)$ exactly: the posterior over
  structures has an exact two-fold symmetry, the modal structure is not unique, and
  \textbf{reporting the mode alone is misleading} --- marginal inclusion
  probabilities split evenly and are the honest summary. Including \emph{both} adds a
  null direction, so
\begin{equation}
\log\frac{\Zc_i(\gamma\cup\{j,j'\},g)}{\Zc_i(\gamma\cup\{j\},g)} \;=\;
-\tfrac12\log(1+g\cdot0) \;+\; 0 \;=\; 0:
\tag{E-129}\label{E129}
\end{equation}
  the likelihood side is \emph{exactly} indifferent, and the only thing separating
  them is the multiplicity term $\approx-\log n_i$. This is the cleanest
  demonstration in the document that the two complexity mechanisms do different jobs
  --- the volume correction charges nothing for a redundant parameter, correctly,
  since it costs no effective dimension; the multiplicity correction charges
  $\log n_i$, also correctly, since a redundant entry is still a claim about the
  structure. A derivation carrying only the first would include duplicates for free.

\item[(D-e) $T\le p_{\max}$.]
  No complete $(y_t,z_t)$ pairs; \eqref{E101} returns the prior predictive. Well
  defined and useless: the program must refuse, and the declaration requires
  $T\ge p_{\max}+2$.

\item[(D-f) A constant coordinate.]
  A coordinate identically zero after differencing gives $c_i=0$, so \eqref{E101}
  attains the bound \eqref{E128}: finite but degenerate. Detect at declaration and
  refuse. \textbf{This is the case AUTOCLASS's per-attribute \texttt{error} absorbs
  and we cannot} (Remark \ref{rem:autoclass}); refusing is the price of taking no
  scale from the user.
\end{description}


\section{Reductions}\label{sec:red}

Charter \S7 makes reductions a gate condition. The first two are anchors: they show
the model class contains the right textbook cases. \textbf{Neither is a correctness
test}, because both are evaluated where every mechanism this project exists to
exercise is switched off. \S\ref{sec:red3} is the one with teeth and it is new.

\subsection{All entries on $\Rightarrow$ the conjugate Bayesian VAR}

\begin{proposition}[reduction 1]\label{prop:red1}
At $\gamma_i=\{0,1,\dots,Np\}$ for every $i$ and conditional on $g$,
\begin{equation}
\Zc(\Sc_{\text{full}},g) = \prod_{i=1}^{N}(2\pi)^{-T/2}\bigl|I+gS_{zz}\bigr|^{-1/2}
\frac{b_\sigma^{a_\sigma}\Gamma(a_\sigma+\frac{T}{2})}{\Gamma(a_\sigma)}
\Bigl(b_\sigma+\tfrac{R_i}{2}\Bigr)^{-(a_\sigma+\frac{T}{2})},
\tag{E-130}\label{E130}
\end{equation}
$R_i=c_i-s_i'(S_{zz}+g^{-1}I)^{-1}s_i$: exactly the conjugate Bayesian VAR with
independent Normal--inverse-Gamma priors per equation, diagonal $\Sigma$, and prior
coefficient covariance $g\sigma_i^2I$ --- the standard closed form, carrying the
$g$-dependence that the standard treatment omits by fixing it.
\end{proposition}

\begin{proof}
Immediate from Proposition \ref{prop:closed}, since at full inclusion \eqref{E29_A}
equals \eqref{E28_A}.
\end{proof}

\subsection{$p=0$ $\Rightarrow$ AUTOCLASS's single-class real model}\label{sec:red2}

At $p=0$, $z_t$ is the constant alone, so every $\gamma_i=\{0\}$ and
\begin{equation}
\Zc(\Sc_{p=0}) = \prod_{i=1}^{N}\iint\prod_t\mathcal{N}(y_{ti};m_i,\sigma_i^2)\,
\mathcal{N}(m_i;0,g\sigma_i^2)\,\operatorname{InvGamma}(\sigma_i^2;a_\sigma,b_\sigma)
\,dm_i\,d\sigma_i^2,
\tag{E-131}\label{E131}
\end{equation}
$N$ independent Normal--inverse-Gamma models, one per coordinate, with a free mean.
With the intercept of \eqref{E134_A} now in the model this is
\begin{equation}
y_{ti}\sim\mathcal{N}(m_i,\sigma_i^2),
\tag{E-132}\label{E132}
\end{equation}
with $m_i$ free, which is \textbf{exactly} AUTOCLASS's single-class real-valued model
(\texttt{single\_normal\_cn}). The previous draft, having no intercept, reduced
instead to that model with its mean pinned at zero, and had to record the mismatch.
One further difference remains and is a repair rather than a discrepancy: AUTOCLASS's
prior on the mean is flat and made normalisable by cutting it off at data-derived
limits, which its own authors call cheating (\texttt{tr-fia-90-12-7-01} \S4.2,
\emph{``in order to make it normalizable we must cheat and use information from the
data''}). Ours is proper and its scale is integrated out.

\subsection[Reduction 3: sixteen structures]{Reduction 3: sixteen structures, brute force, everything at
once}\label{sec:red3}

\begin{recordbox}
\textbf{Why this exists.} Reduction 1 holds $g$ fixed at full inclusion: the
structure prior is a single number compared with nothing, $\Delta_\pi$ is never
exercised, no structure move happens, the quadrature is bypassed, and the Occam
factor is evaluated at one point. Reduction 2 has no $g$, no $\gamma$ and no
quadrature at all. Between them they test \eqref{E101} at a single structure --- the
same object the $N=1,k=1,T=2$ fixed-dimension check already tests. \textbf{They
cannot catch} a sign error in \eqref{E113}, the wrong $n_i$ in the multiplicity term,
a quadrature that truncates a structure-dependent tail, or the $(k+1)$ error of
\S\ref{sec:inclusion}. This reduction can catch all four.
\end{recordbox}

\begin{definition}[the check]\label{def:check}
Take $N=2$, $p=1$, $T=4$, fixed rational $y_0,\dots,y_4\in\mathbb{R}^2$ (five
vectors: four usable transitions and one pre-sample). Each equation has $n_i=Np=2$
candidates plus the always-included intercept, so the structure space has
$2^n=2^4=16$ members. Compute the full posterior
\begin{equation}
\pi(\Sc\mid y) \;=\; \frac{\pi(\gamma)\,\Zc(\Sc)}{\sum_{\Sc'}\pi(\gamma')\,\Zc(\Sc')},
\tag{E-136}\label{E136}
\end{equation}
with $\Sc$ ranging over all 16, \textbf{two ways}, and require agreement to $10^{-6}$ in each of the sixteen
probabilities and to $10^{-4}$ nats in each of the sixteen scores.
\begin{description}[style=nextline, leftmargin=2em, topsep=0pt, itemsep=0pt]
\item[Reference.]
  For each structure, integrate the joint \eqref{E98} numerically over
  $(\alpha_1[\gamma_1],\sigma_1^2,\alpha_2[\gamma_2],\sigma_2^2)$ and over $g$ --- at
  most three dimensions per equation, so tensor-product Gauss quadrature reaches
  machine precision cheaply --- \emph{without} using \eqref{E101}. Multiply by
  $\pi(\gamma)$ computed directly from \eqref{E96_A}--\eqref{E96_B} as
  $\prod_i\int_0^1\rho_i^{k_i}(1-\rho_i)^{n_i-k_i}d\rho_i$, i.e.\ without using
  \eqref{E111_A} or \eqref{E113}. Normalise over the sixteen.
\item[Production path.]
  \eqref{E101}$\to$\eqref{E102}$\to$ \eqref{E103}/\eqref{E104} on the shipped grid
  $\to$\eqref{E135}$\to$ \eqref{E126}$\to$\eqref{E136}.
\end{description}
\textbf{The reference must use a different quadrature rule, a different node count
and a different order of integration from the production path}, so that a shared bug
cannot cancel.
\end{definition}

\paragraph{What it exercises, and what a failure would look like.} It exercises
\eqref{E98}, \eqref{E101}, \eqref{E111_A}, \eqref{E113}, \eqref{E135}, the
$g$-quadrature and the normalisation \eqref{E136} \emph{simultaneously}, and it fails
if any one of them is wrong. The failure signatures are diagnostic: a sign error in
\eqref{E113} shows as a monotone drift in the discrepancy with $k=\sum_ik_i$; a wrong
$n_i$ (that is, $N^2p$ where $Np$ belongs, or a candidate count that wrongly includes
the intercept) shows as a constant offset per included entry; a truncated quadrature
tail shows as a discrepancy growing with $k$, because by \eqref{E107_B} the right tail
slope depends on $\sum_ik_i^{\eff}$; and an error in \eqref{E101} itself shows at
$k=0$, where everything else is inert.

$T=4$ is chosen so that $k_i\le2<T$ and $S_\gamma$ is non-singular with a non-trivial
fit term, while the whole computation remains exactly reproducible from five printed
rational vectors. Repeating at $T=3$ and $T=6$ costs nothing and catches a
$T$-indexing error.


\section{What the model class cannot represent}\label{sec:cannot}

This is the one place v1 is provably unable to represent something the data may
contain, and it is a consequence of fixing $d=1$ without inference.

\begin{proposition}[the observed spectral density is positive definite at every
frequency]\label{prop:spec}
Write \eqref{E1_A} as $\mathcal{A}(L)y_t=\varepsilon_t$ with
\begin{eqnarray}
\mathcal{A}(L) &=& I_N - \sum_{l=1}^{p}A_lL^l,
  \Elabel{E-6A}{E6_A} \\
\mathcal{A}(0) &=& I_N,
  \Elabel{E-6B}{E6_B}
\end{eqnarray}
and let $\theta$ satisfy $\det\mathcal{A}(e^{-i\omega})\ne0$ for all $\omega$. Then
\begin{equation}
f_y(\omega) \;=\; \tfrac{1}{2\pi}\,\mathcal{A}(e^{-i\omega})^{-1}\,\Sigma\,
\bigl(\mathcal{A}(e^{-i\omega})^{-1}\bigr)^{*} \;\succ\; 0
\tag{E-76}\label{E76}
\end{equation}
for every $\omega$, every $\theta$, every $p$.
\end{proposition}

\begin{proof}
\eqref{E76} is the filtering identity. For definiteness let $v\ne0$ and
$w=(\mathcal{A}(e^{-i\omega})^{-1})^*v$; then $2\pi v^*f_yv=w^*\Sigma w$, and $w\ne0$
by invertibility while $\Sigma\succ0$ by \eqref{E4_B}.
\end{proof}

\begin{corollary}[over-differenced data lies outside the model class]\label{cor:overdiff}
A spectral zero --- a frequency at which the true spectral density is singular,
equivalently a unit root in the moving-average operator --- is produced by no
$\theta\in\Theta$, at any $p$. With no hidden state there is not even a channel
through which the misfit could be reported as something else: it is absorbed by $p$
and by the sparsity pattern, which are the two quantities the project exists to
report.
\end{corollary}

\paragraph{What the machinery does instead, derived.} Suppose a coordinate is
stationary in levels and the pipeline differences it. In the pure case,
$y_t=e_t-e_{t-1}$ is MA(1) with coefficient $-1$; its autocovariances are
\begin{eqnarray}
\gamma_0 &=& 2\varsigma^2,
  \Elabel{E-77A}{E77_A} \\
\gamma_1 &=& -\varsigma^2,
  \Elabel{E-77B}{E77_B} \\
\gamma_r &=& 0 \quad (r\ge2),
  \Elabel{E-77C}{E77_C}
\end{eqnarray}
and the Yule--Walker equations become a discrete Laplacian with fixed ends, whose
solution is linear in the lag index:
\begin{equation}
\varphi_m^{(p)} \;=\; -\Bigl(1-\frac{m}{p+1}\Bigr),
\tag{E-78}\label{E78}
\end{equation}
for $m=1,\dots,p$, and
\begin{eqnarray}
v_p &=& \varsigma^2\,\frac{p+2}{p+1},
  \Elabel{E-79A}{E79_A} \\
v_1 &=& \tfrac32\varsigma^2,
  \Elabel{E-79B}{E79_B} \\
v_2 &=& \tfrac43\varsigma^2,
  \Elabel{E-79C}{E79_C} \\
v_\infty &=& \varsigma^2.
  \Elabel{E-79D}{E79_D}
\end{eqnarray}
Three readings. (i) \textbf{The fit improves monotonically in $p$ and never attains
its limit}: there is no $p$ at which the data stops asking for another lag. (ii) At
any fixed $p$ there \emph{is} an interior maximiser, namely \eqref{E78}, so the
pathology is not that a fit fails to converge but that the \emph{sequence} of fits
does not terminate. (iii) \textbf{Therefore the posterior over $p$ on
over-differenced data is decided by $\pi(p)$ and by $p_{\max}$, not by the data}, and
the spurious links land on the diagonal --- where the analyst is least likely to
question them.

\paragraph{Mitigation, and its honest limit.} Open item O-5 proposes fitting at both
$d=0$ and $d=1$ and presenting both. The derivation supports the proposal and
sharpens the caveat: the two fits are models for \emph{different data}, so putting
them side by side is legitimate and subtracting them is not. \textbf{The choice rests
on the analyst and on diagnostics, not on the machinery}, and this document declines
to supply a number that would let anyone pretend otherwise. What \emph{is} comparable
is a within-$d$ diagnostic:
\begin{equation}
\kappa(p) \;:=\; \Bigl|\det\Bigl(I_N-\sum_{l=1}^{p}\hat A_l\Bigr)\Bigr|,
\tag{E-80}\label{E80}
\end{equation}
with over-differencing showing as $\kappa(p)$ small and decreasing in $p$, which
\eqref{E78} predicts exactly, since $1-\sum_m\varphi_m^{(p)}z^m$ at $z=1$ is
$1/(p+1)\to0$. It costs nothing, it is a report and not a score, so D-10 is
untouched. If a coordinate trips it, the honest statement is that $d=1$ was wrong for
that coordinate and its equation's reported structure is not interpretable.

\begin{recordbox}
\textbf{Self-flagged weakness (W-2, rank 2).} Corollary \ref{cor:overdiff} is
airtight; the \emph{quantification} is worst-case only. \eqref{E77_A}--\eqref{E79_D} are
computed at the extreme (white in levels); a coordinate that is nearly a unit root in
levels is barely distorted by differencing, and nothing in v1 tells the analyst which
they have. \eqref{E80} detects the extreme end only. \textbf{This is the largest
known hole in v1}, and it is a scope consequence of D-05 rather than an error. O-5
needs a ruling before P2.
\end{recordbox}


\section{The conditioning diagnostic}\label{sec:cond}

Near-flat directions arise here from \textbf{collinear regressors}, which is live in
the target regime: at $N=50$, $p=2$ there are $n_i=100$ candidates per equation
against 198 usable transitions, and the columns are lagged values of a strongly
cross-correlated process, not independent draws. The object is
\begin{eqnarray}
H_{\ell,i}[\gamma_i,\gamma_i] &=& \tfrac{1}{\sigma_i^2}S_\gamma,
  \Elabel{E-82A}{E82_A} \\
H_i[\gamma_i,\gamma_i] &=& \tfrac{1}{\sigma_i^2}\bigl(S_\gamma+\tfrac1gI_k\bigr),
  \Elabel{E-82B}{E82_B}
\end{eqnarray}
so the two differ by exactly the slab precision and the spectrum of $S_\gamma$
carries everything.

\begin{proposition}[what a rank computation at fixed $T$ cannot see]\label{prop:rank}
Let $v$ be a fixed unit vector. (a) If $\E[(v'z_t)^2]=\varrho^2>0$ then
$v'S_\gamma v/T\to\varrho^2$ almost surely: the direction is data-informed and its
eigenvalue grows linearly in $T$. (b) If $\varrho^2=0$, exact collinearity, then
$v'S_\gamma v=0$ for every $T$ and $v'H_iv=1/(g\sigma_i^2)$, independent of $T$: the
direction is carried entirely by the prior. (c) If $\varrho^2$ is small but non-zero
--- the realistic case --- the direction is data-informed but becomes visible only
once $T\varrho^2$ exceeds $1/g$, and \textbf{at fixed $T$ it is indistinguishable from
(b) by any rank computation at any tolerance}, since the classification depends only
on the product $T\varrho^2$.
\end{proposition}

\begin{proof}
(a),(b) are the law of large numbers applied to $v'S_\gamma v=\sum_t(v'z_t)^2$. For
(c), a tolerance $\varsigma$ classifies by $\mathbf{1}[T\varrho^2>\varsigma]$, which
is a statement about the product; two processes with
$T_1\varrho_1^2=T_2\varrho_2^2$ are classified alike while differing in whether the
direction is identified in the limit.
\end{proof}

\begin{definition}[the diagnostic --- a report, never a score]\label{def:diag}
Across the sweep, report the decade histogram of the spectrum,
\begin{equation}
\operatorname{bucket}(\mu_j) = \bigl\lfloor\log_{10}(\mu_j/\mu_1)\bigr\rfloor,
\tag{E-83}\label{E83}
\end{equation}
and, across two sample sizes on the same process, the scaling exponent
\begin{equation}
r_j(T_1,T_2) = \frac{\log(\mu_j(T_2)/\mu_j(T_1))}{\log(T_2/T_1)},
\tag{E-84}\label{E84}
\end{equation}
which tends to 1 for a data-informed direction and 0 for a prior-carried one.
\textbf{Never report a tolerance-based rank.} Neither quantity enters any score.
\end{definition}

Under independent regressors at $N=50$, $p=2$, $T=200$ the condition number of
$S_{zz}$ would be near 35, which is comfortable; under B1's actual generator, whose
within-group blocks are dense with $\rho=0.7$, the small end of the spectrum is not
protected by that arithmetic, and \eqref{E84} across the sweep is the only thing that
will say which regime the data is in.


\section[The observation model]{The observation model as a declared component}\label{sec:obs}

AUTOCLASS's engine character came from a library of likelihood terms selected per
attribute --- a multinomial term, a normal term, a normal-with-modelled-missingness
term, a covariant block --- with its model file doing nothing but assigning terms to
attribute indices. This section states, honestly and briefly, how far that
transplants.

The derivation is carried for a general conjugate exponential-family regression: a
declared type supplies a base measure, a sufficient-statistic map $\mathsf{s}_i$, a
log-partition $B_i$ and a dispersion $\tau_i$, with the autoregression carrying the
natural parameter through a link,
\begin{equation}
f(y_{ti}\mid\nu_{ti},\tau_i) = h_i(y_{ti};\tau_i)\exp\bigl\{\tfrac{1}{\tau_i}
[\nu_{ti}\mathsf{s}_i(y_{ti})-B_i(\nu_{ti})]\bigr\},
\tag{E-87}\label{E87}
\end{equation}
\begin{equation}
\nu_{ti} \;=\; \alpha_i[\gamma_i]'z_t[\gamma_i] \qquad \text{(canonical link)},
\tag{E-88}\label{E88}
\end{equation}
\begin{eqnarray}
\ell_i &=& \sum_t\log h_i + \tfrac{1}{\tau_i}\Bigl[\alpha_i[\gamma_i]'\mathsf{s}_i^z
  - \sum_tB_i\bigl(\alpha_i[\gamma_i]'z_t[\gamma_i]\bigr)\Bigr],
  \Elabel{E-89A}{E89_A} \\
\mathsf{s}_i^z &:=& \sum_t\mathsf{s}_i(y_{ti})z_t[\gamma_i].
  \Elabel{E-89B}{E89_B}
\end{eqnarray}
Taking $\mathsf{s}_i(y)=y$, $B_i(\nu)=\nu^2/2$, $\tau_i=\sigma_i^2$ recovers
\eqref{E29_A} exactly:
\begin{equation}
\ell_i = -\tfrac{T}{2}\log(2\pi\sigma_i^2) - \tfrac{1}{2\sigma_i^2}\bigl(c_i
- 2\alpha_i[\gamma_i]'s_i[\gamma_i] + \alpha_i[\gamma_i]'S_\gamma\alpha_i[\gamma_i]\bigr),
\tag{E-90}\label{E90}
\end{equation}
so the Gaussian case is an instantiation and not a parallel derivation. A type is
admissible if it supplies (T1) the four objects of \eqref{E87}--\eqref{E88}; (T2)
statistics of fixed dimension, independent of $T$, accumulable in one pass; (T3) a
proper prior on $(\alpha_i,\tau_i)$ into which the slab scale enters as in
\S\ref{sec:prior}; and (T4) an evaluator for
\begin{equation}
\Zc_i(\gamma_i,g) = \iint e^{\ell_i}\,\pi(\alpha_i[\gamma_i]\mid\tau_i,g)\,
\pi(\tau_i)\,d\alpha_i\,d\tau_i.
\tag{E-91}\label{E91}
\end{equation}

\begin{recordbox}
\textbf{What this does and does not cover, said plainly.}

\emph{Does.} Everything in \S\ref{sec:prior}, \S\ref{sec:inclusion},
\S\ref{sec:lag}, \S\ref{sec:search} and \S\ref{sec:red} is written in terms of
$\log\Zc_i(\gamma_i,g)$ and $|\gamma_i|$ alone and \textbf{knows nothing about the
coordinate type}. That is the precise and entire sense in which the sparsity and
complexity machinery is type-indifferent.

\emph{Does not.} \textbf{One type is derived: Gaussian. No other type is derived, and
(T2) and (T3) are asserted per type and verified for none.} A count coordinate under
a log-link Poisson autoregression has $B_i(\nu)=e^\nu$, so the second term of
\eqref{E89_A} is not quadratic, no finite-dimensional conjugate prior on $\alpha_i$
exists, and \eqref{E91} has no closed form; the same is true of a multinomial
coordinate under a logit link. The price is a Laplace evaluation or a
low-dimensional quadrature per equation per candidate structure --- the \emph{same
species} of price D-24 already accepts for the slab scale, which is the honest
framing, but it is not free. Incremental updating degrades from $O(k^2)$ to an
iterative solve, since Proposition \ref{prop:flip} rests on the bordered Cholesky of
a matrix that no longer stays fixed. And a stationarity notion may not exist at all
for a non-Gaussian autoregression, which \S\ref{sec:cond} and \S\ref{sec:cannot}
both assume.

\emph{In AUTOCLASS the analogue was nearly free} because attributes are independent
given the class, so terms of different types simply multiply. Here the analogue of an
attribute term is a \emph{regression}, and a regression needs a link. Claiming the
transplant is complete would be overselling it.
\end{recordbox}


\section{Assumptions and limitations}\label{sec:assump}

\paragraph{Assumptions in force.} \textbf{A1} Conditioning on $y_{1-p_{\max}:0}$ with
$p_{\max}$ declared before fitting (Assumption \ref{ass:cond}). \textbf{A2} Every
prior proper (Assumption \ref{ass:proper}). \textbf{A3} No stationarity truncation,
so the prior places mass on explosive processes (Remark \ref{rem:trunc}).
\textbf{A4} $\Sigma$ diagonal (D-04), consumed at \eqref{E28_A}; contemporaneous
innovation correlation is \textbf{not representable in v1 at all}, and with $K=0$
there is no hidden coordinate to absorb it. \textbf{A5} $d$ fixed at 1, not inferred,
not comparable across (D-05). \textbf{A6} Coordinates standardized before fitting, on
which \eqref{E92_B} and \eqref{E95_A} are calibrated --- so the prior is empirical-Bayes
to the extent of one scalar per coordinate. It cancels from every Bayes factor and
does not cancel from a reported $\log\Zc$. \textbf{A7} Regular sampling, no missing
values. \textbf{A8} $A_l$ constant in $t$; no regime switching, no time variation.
\textbf{A9} $z_t$ stationary and ergodic, used in Proposition \ref{prop:rank}(a) and
in \eqref{E124_A}; not implied by A3.

\paragraph{Limitations.} \textbf{L1} A pure VAR cannot represent moving-average
structure at all. Any genuine MA component must be approximated autoregressively, and
\eqref{E79_A} shows how badly: the approximation improves like $(p+2)/(p+1)$ and never
terminates. \textbf{L2} Over-differenced data lies outside the model class
(\S\ref{sec:cannot}); mitigation is \eqref{E80} plus O-5, which is \textbf{open and
needs a ruling before P2}. \textbf{L3} No hidden state, hence no latent common
factor: a common shock is representable only as explicit links among observed
coordinates, at up to $N$ entries per shock rather than one dimension. \textbf{L4}
Missing values are unsupported; the charter says otherwise and should be amended. A
missing component removes one observation from one equation and one from every
equation at the next step, breaking the equal-effective-sample condition
\emph{across} equations. \textbf{L5} Real-valued coordinates only; no other type is
derived (\S\ref{sec:obs}). \textbf{L6 No inferred grouping is reported.} A reader who
wants a partition of the coordinates must read it off the recovered sparsity pattern.
This is a deliberate consequence of D-30 and is arguably the better behaviour --- we
report what the data supports and do not impose the concept of a group at all --- but
it is a capability the previous draft claimed and this one does not. \textbf{L7} The
reported posterior is normalised over the retained set, not over all structures
(W-4). \textbf{L8} Cross-$d$ and cross-$N$ comparisons are meaningless; only $p$ and
structure are compared. \textbf{L9} With a diagonal $\Sigma$ assumed, a structure
that happens to leave less contemporaneous correlation unexplained will score better
for that reason and not because its link structure is truer. AUTOCLASS records the
identical distortion from its own independence assumption
(\texttt{interpretation-c.text} 30--37) and it belongs in the limitations in those
terms. Benchmark variant B1-N measures it; \textbf{the size is not derived here}
(W-3).


\appendix

\section{Equation disposition}\label{app:disp}

Identifiers are permanent. A withdrawn or retired equation keeps its identifier and
is never reused. New equations take the next free number regardless of where they
sit, so the table is in identifier order.

{\small
\begin{longtable}{@{}>{\raggedright\arraybackslash}p{4.4cm}
                    >{\raggedright\arraybackslash}p{3.4cm}
                    >{\raggedright\arraybackslash}p{8.9cm}@{}}
\toprule
IDs & Disposition & Note \\
\midrule
\endfirsthead
\toprule
IDs & Disposition & Note \\
\midrule
\endhead
\bottomrule
\endlastfoot
E-1, E-2, E-4--E-6, E-9 & restated &
  $x=y$, $J=I_N$; $D=N^2p+2N$ now carries the intercept \\
E-3, E-7, E-8, E-11--E-13, E-31 & withdrawn (round 2) &
  hidden-state bookkeeping; dormant \\
E-10 & unchanged & complete-data likelihood, folded into \eqref{E9} \\
E-14--E-17 & unchanged &
  joint, $\Zc$, structure posterior, the exact identity; cited not reprinted \\
E-18--E-20 & unchanged &
  Laplace, Occam, BIC; cited in Remark \ref{rem:laplace}, not used \\
E-21 & restated & $D_{\eff}=D$ at group dimension 0 \\
E-22, E-23 & restated (round 2) &
  missing information $\equiv0$; trivial, not reprinted \\
E-24--E-30, E-32 & unchanged &
  score; regression form; sufficiency; decoupling; bordering \\
E-33--E-75 & \textbf{withdrawn, dormant} &
  hidden state (A-001); live in \texttt{theory-v0-hidden-state.tex} \\
\midrule
E-76--E-86 & unchanged &
  spectral density, over-differencing, Hessian conventions, diagnostic, flip cost \\
E-87--E-91 & unchanged, compressed &
  exponential-family interface, \S\ref{sec:obs} \\
E-92--E-97 & unchanged &
  hyperpriors; every constant now also in \S\ref{sec:constants} \\
E-98--E-110 & unchanged &
  the closed form, the quadrature, the cache, the candidate set \\
E-111--E-113 & \textbf{restated} &
  per-equation counts $(k_i,n_i)$ in place of the global $(k,n)$ ---
  \S\ref{sec:rate}, W-1. Formulae identical; indexing changed \\
\textbf{E-114--E-117} & \textbf{RETIRED (D-30)} &
  model S-B: induced entry pattern, Ewens/CRP partition prior, $\alpha_{\text{CRP}}$
  hyperprior, block indicators. \textbf{Never reused.} Live in
  \texttt{theory-v1-with-blocks.tex} \\
E-118, E-119, E-121 & unchanged &
  partial $t$-statistic; three-term increment; sweep prediction \\
\textbf{E-120, E-122, E-133} & \textbf{corrected (O-004)} &
  the $(k+1)$ term restored; every dependent number in \S\ref{sec:predict}
  recomputed \\
E-123--E-125 & restated &
  $N=2$, $p=2$ design (intercept added); $p^*$ demoted to advisory (O-003) \\
E-126 & restated & the score, with the S-B branch deleted \\
E-127--E-132 & restated &
  degenerate cases and reductions 1--2; \eqref{E132} now an exact match rather than
  a discrepancy, the intercept having closed it \\
\midrule
\textbf{E-134} & \textbf{new} & regressor vector with always-included intercept \\
\textbf{E-135} & \textbf{new} & joint structure prior as a product over equations \\
\textbf{E-136} & \textbf{new} &
  reduction 3: the 16-structure exact posterior, \S\ref{sec:red3} \\
\textbf{E-137} & \textbf{new} &
  the declared evidence threshold and the chain hanging from it,
  \S\ref{sec:constants} \\
\end{longtable}
}

\noindent\textbf{Agent 3's work list, in priority order.}
\begin{enumerate}[topsep=0pt, itemsep=0pt]
\item \textbf{Reduction 3} (Definition \ref{def:check}). It subsumes items 2 and 4 of
  the previous list and is the only check that fails if any single component is
  wrong.
\item \eqref{E101} at $N=1$, $k=1$, $T=2$ against the direct double integral
  \eqref{E98}. Cheap, and it isolates the one object everything rests on.
\item The $N=2$, $p=2$ battery of \S\ref{sec:inst}. \textbf{No $p=2$ check exists
  anywhere in the project.}
\item \eqref{E94} --- the Lindley--Bartlett collapse, owed since round 1 and never
  reached.
\item \eqref{E119} at $k=0\to1$ and $1\to2$ \textbf{with the corrected}
  \eqref{E120}; then \eqref{E129}.
\item \eqref{E130}, \eqref{E131}, the two anchor reductions.
\item \eqref{E21_B}/\eqref{E81_C}: $\operatorname{rank}H_\ell=\operatorname{rank}H=D$ at
  $N=2$, $p=1$, $T=3$.
\item Numeric: \eqref{E84} on a constructed collinear design; then \eqref{E108}.
\end{enumerate}
Every round-1 script for E-33--E-75 is dormant; every script written against the
\emph{global}-rate forms of E-111--E-113 must be re-run with per-equation counts.

\section{Self-flagged weaknesses, ranked}

\begin{description}[leftmargin=3em, labelindent=0pt, topsep=0pt, itemsep=0pt]
\item[W-1.] \textbf{Per-equation versus global inclusion rate} (\S\ref{sec:rate}).
  Worth $\log N=3.9$ nats per entry at $N=50$ and a factor of $N$ in the predicted
  false-positive count. Both columns are printed in Table \ref{tab:b1}.
  \textbf{Confirmation requested.}
\item[W-2.] \textbf{Over-differencing is a model-class failure and v1 detects it only
  at the extreme} (\S\ref{sec:cannot}). Airtight as a statement; quantified only at
  the worst case. \textbf{O-5 needs a ruling before P2.}
\item[W-3.] \textbf{The misspecified-$\Sigma$ distortion is not quantified} (L9). The
  mechanism is derived and the size is not; B1-N measures it.
\item[W-4.] \textbf{The reported posterior is normalised over the retained set}
  (\S\ref{sec:search}). Mitigation offered is a report, not a bound.
\item[W-5.] \textbf{$\delta_{\text{report}}=1$ nat is declared, not derived}
  (\S\ref{sec:constants}). The argument constrains it to be well below $\log n_i$ and
  well above the quadrature error, which is a range, not a number. A different
  declaration changes which near-threshold entries are called decided.
\item[W-6.] \textbf{$\varepsilon$ in \eqref{E124_A} is a reporting convention I cannot
  derive}, and at $\varepsilon=\frac14$ the charter's own target does not clear $p=2$
  (\S\ref{sec:ceiling}). I decline to pick the value that makes the charter right.
\item[W-7.] \textbf{Standardization makes \S\ref{sec:prior} empirical-Bayes by one
  scalar per coordinate} (A6). Cancels from every Bayes factor; does not cancel from
  a reported $\log\Zc$.
\item[W-8.] \textbf{Non-Gaussian coordinate types are specified, not derived}
  (\S\ref{sec:obs}). One instantiation is worked; (T2)/(T3) are verified for no other
  type.
\item[W-9.] \textbf{Missing values are unsupported and the charter says otherwise}
  (L4). Stated, not solved.
\end{description}

\section{What was deleted, and where it lives}

\textbf{By D-30 (this revision).} Model S-B in full: the partition $\mathcal{P}$ of
the coordinates into $G$ groups, the Ewens/CRP prior over partitions with its
concentration hyperprior, the block-level inclusion indicators $\delta_{gh}^{(l)}$
and the induced entry pattern, the group-count inference, the S-A/S-B margin
arithmetic, the prior over the pair of models, the split/merge/move search operators,
the two declaration fields that selected between the models, and every predicted sign
of a comparison between them. Equations \textbf{E-114--E-117}, retired. All of it is
preserved verbatim in \texttt{theory/theory-v1-with-blocks.tex}, which carries a
boxed supersession notice.

\textbf{By A-001 (earlier).} The hidden state and everything attached to it: the
$GL(K)$ symmetry, both identifiability routes, the Faddeev--Popov determinant, orbit
volumes, pseudo-determinant Laplace, the prior over $K$, the hidden pre-sample
initialisation and the $K$-dependent constant it put in the score, the
Cheeseman--Stutz debt, and the EM/Kalman inner loop. Equations \textbf{E-33--E-75},
dormant in \texttt{theory/theory-v0-hidden-state.tex}. Three confirmed blocking
defects against that material (V-001--V-003) are unrepaired and must be the first
items addressed if Extension 2 resumes.

\end{document}

% ---------------------------------------------
% Copyright Ben Paul Wise. All rights reserved.
% ---------------------------------------------
